\documentclass[journal,onecolumn,11pt]{IEEEtran}
% NOTE: IEEE-стилско форматирање (IEEEtran). Компилирано со XeLaTeX поради
% поддршка на кирилица; распоредот/форматирањето е идентично со оригиналот,
% само фонтот е заменет со FreeSerif (Times-сличен, со полна кирилична поддршка).
% Овој документ ги обединува двата извештаи за истиот BLDC/PMSM инвертор:
%   (1) енергетскиот степен и колото за управување на гејтот;
%   (2) аналогниот влезен степен (мерење на струја/напон, snubber, филтрирање).

\usepackage{fontspec}
\setmainfont{FreeSerif}
\setsansfont{FreeSans}
\setmonofont{FreeMono}
\usepackage{amsmath,amssymb}
\usepackage{mathtools}
\usepackage{babel}
\babelprovide[import, main]{macedonian}
\usepackage{siunitx}
\usepackage{booktabs}
\usepackage{array}
\usepackage{tabularx}
\usepackage{ragged2e}
\usepackage{multirow}
\usepackage{graphicx}
\usepackage[siunitx]{circuitikz}
\usepackage{tikz}
\usetikzlibrary{arrows.meta,positioning,calc,shapes.geometric}
\usepackage[hidelinks,colorlinks=true,linkcolor=blue,citecolor=blue,urlcolor=blue]{hyperref}
\usepackage{url}
\usepackage{xcolor}
\usepackage{enumitem}
\usepackage{lettrine}
% IEEEtran's native \IEEEPARstart drop-cap uses TFM font loading that is
% incompatible with XeLaTeX/fontspec; redefine it with lettrine instead.
\renewcommand{\IEEEPARstart}[2]{\lettrine[lines=2,findent=2pt,nindent=0pt]{#1}{#2}}

% Кратенка за термини задржани на англиски (нема еднозначен превод или преводот
% би го изгубил значењето). Дефинирана како поминувачка макро за да биде
% форматирањето идентично со останатиот текст.
\newcommand{\eng}[1]{#1}

\sisetup{detect-all,per-mode=symbol,group-digits=integer,range-phrase={ до }}
\DeclareSIUnit{\ohm}{\text{Ω}}

% Тип на колона со порамнување налево и автоматско прекршување (за широки табели)
\newcolumntype{L}{>{\RaggedRight\arraybackslash}X}

% Македонски имиња на структурните ознаки (IEEE-стил)
\AtBeginDocument{%
  \def\figurename{Сл.}%
  \def\tablename{ТАБЕЛА}%
  \def\refname{Референци}%
  \def\abstractname{Апстракт}%
  \def\IEEEkeywordsname{Клучни зборови}%
}
% потолерантно прекршување на редови за неделив кириличен текст
\tolerance=2000
\emergencystretch=2em
\hbadness=6000

% Convenience macros
\newcommand{\Vdcbus}{V_{\mathrm{bus}}}
\newcommand{\Rshunt}{R_{\mathrm{sh}}}
\newcommand{\Iphrms}{I_{\mathrm{ph}}}
\newcommand{\Fswsw}{f_{\mathrm{sw}}}

\begin{document}

\title{Дизајн на трифазен инвертор со N-канални MOSFET транзистори за
\mbox{BLDC/PMSM} мотор од \mbox{3--\SI{5}{\kilo\watt}}, напојуван од
\mbox{24-V/48-V} батерии (\mbox{20--\SI{60}{\volt}}): енергетски степен,
управување на гејтот и аналоген влезен степен со векторско управување (FOC)}

\author{Извештај за хардверски дизајн\\
\small Целен контролер: STMicroelectronics STM32G473RCT6 \quad
Енергетски степен: Infineon IPT015N10N5 (\SI{100}{\volt}, \SI{1.5}{\milli\ohm}),
$\times 2$ паралелно по позиција \quad Драјвер: Texas Instruments UCC27211A-Q1\\
\small Склопна фреквенција \SI{40}{\kilo\hertz} \quad Управување на гејтот:
\SI{12}{\volt} \quad Струен засилувач: TI INA240A2 \quad Набавка: LCSC Electronics}

\markboth{Извештај за дизајн на трифазен BLDC/PMSM инвертор}{}
\maketitle

\section{Вовед и опсег}
\IEEEPARstart{Ц}{елниот} систем е неизолиран трифазен напонски инвертор со
заедничко заземјување (common-ground), наменет да управува со произволен трифазен
безчеткест еднонасочен или синхрон мотор со перманентни магнети (BLDC/PMSM), под
сензорско управување (Hall или квадратурен енкодер) или безсензорско, базирано на
повратна електромоторна сила (back-EMF), со целосно векторско управување (FOC).
Системот се напојува од батерија со променлив напон: пакет од \SI{24}{\volt}
(опсег \mbox{20--\SI{30}{\volt}}), кој дава околу \SI{3}{\kilo\watt}, или пакет од
\SI{48}{\volt} (опсег \mbox{40--\SI{60}{\volt}}), кој дава до \SI{5}{\kilo\watt}.
Промената на напонот на батеријата е суштинска за дизајнот: \emph{највисокиот}
напон (\SI{60}{\volt}) ги одредува напрезите на блокирање и брзината на нишање
$\mathrm{d}v/\mathrm{d}t$, додека \emph{најнискиот} напон (\SI{20}{\volt}) ја дава
најголемата струја за дадена моќност и оттаму најлошиот случај и за спроводните
загуби и за мерењето на струја. Склопната фреквенција е \SI{40}{\kilo\hertz}, над
границата на човечкиот слух, така што PWM тонот е нечуен.

Контролерот е STM32G473RCT6 --- motor-control MCU базиран на
Arm\textsuperscript{\textregistered} Cortex\textsuperscript{\textregistered}-M4,
кој интегрира пет \mbox{12-битни} ADC-а, шест операционни засилувачи употребливи во
режим со програмабилно засилување (PGA) или како follower, седум аналогни
компаратори, три напредни тајмери за управување со мотор и бафер за напонска
референца (VREFBUF)~\cite{stm32g473ds}. Секоја прекинувачка позиција во енергетскиот
степен е составена од \emph{два} паралелно поврзани Infineon
OptiMOS\textsuperscript{\texttrademark}~5 IPT015N10N5 транзистори (\SI{100}{\volt},
\SI{1.5}{\milli\ohm}, \SI{300}{\ampere})~\cite{ipt015ds,ipt015part} --- четири уреди
по гранка на полумостот, дванаесет вкупно --- а секоја гранка ја управува по еден
gate driver Texas Instruments UCC27211A-Q1, напојуван од
\SI{12}{\volt}~\cite{ucc27211}. Фазната струја се мери со in-line shunt, читан од
струен засилувач TI INA240A2~\cite{ina240ds}.

Извештајот го опфаќа целиот хардверски синџир од батеријата до ADC-то на
контролерот и е организиран во следните проектни блокови:
\begin{enumerate}[leftmargin=1.4em]
  \item \emph{избор на енергетскиот MOSFET и проверка на електричната
        компатибилност} со драјверот за целиот опсег на напони на батеријата
        (Дел~\ref{sec:power});
  \item \emph{колото за управување на гејтот} --- bootstrap кондензатор и диода,
        отпорници на гејтот, асиметрична мрежа за вклучување/исклучување,
        pull-down отпорници и кондензатори за декуплирање (Дел~\ref{sec:gatedrive});
  \item ланецот за \emph{мерење на фазна струја} (in-line shunt $+$ засилувач $+$
        ADC интерфејс), за сите три фази (Дел~\ref{sec:current});
  \item ланците за \emph{мерење на напон} на DC-сабирницата и на секој од трите
        фазни јазли (Дел~\ref{sec:voltage});
  \item \emph{RC snubber} паралелно со енергетските прекинувачи
        (Дел~\ref{sec:snubber}); и
  \item \emph{стратегијата за шум и филтрирање} (Дел~\ref{sec:noise}), консолидирана
        со целосна спецификација на материјали (Дел~\ref{sec:bom}).
\end{enumerate}
За секој блок извештајот (а) го објаснува принципот на работа и потребните
компоненти, (б) ги изведува управувачките равенки, (в) пресметува нумерички
вредности за тековната спецификација, (г) избира конкретна компонента достапна од
дистрибутер и (д) го оправдува изборот. Поглавјето со теоријата на работата
(Дел~\ref{sec:theory}) ги поставува концептуалните основи --- зошто е потребен
инвертор за безчеткеста машина, како работат полумостот и bootstrap напојувањето и
која е улогата на \eng{Miller plateau}-то --- пред да се премине на пресметките на
ниво на компоненти. Две посебни прашања се обработени детално: споредбата на
топологиите за мерење на струја што го мотивира in-line изборот
(Дел~\ref{sec:csurvey}), и потребата (или непотребата) од op-amp бафер пред ADC
(Дел~\ref{sec:buffer}).

Секоја квантитативна тврдба и секоја проектна равенка е поткрепена со примарен
извор --- каталошки лист на производителот или призната апликативна белешка --- со
активни линкови. Дополнително, секоја избрана компонента мора да биде нарачлива од
дистрибутерот LCSC Electronics, при што се претпочитаат компоненти со голем обрт за
да биде реална нарачка со количини од \(500+\); LCSC држи на залиха повеќе од
\num{560000} SKU од \(2600+\) производители~\cite{lcschome}. Бидејќи залихата кај
дистрибутерот се менува секојдневно, спецификацијата на материјали во
Дел~\ref{sec:bom} ја означува единствената специјална компонента (shunt-от со
голема моќност) за потврда на залиха при пуштање на нарачката и нуди алтернатива
составена од крајно вообичаени компоненти.

\section{Теорија на работата}\label{sec:theory}

\subsection{Зошто се користи инвертор за погон на BLDC мотор}
Безчеткестиот еднонасочен мотор е, и покрај своето име, синхрона машина за
наизменична струја: има ротор со траен магнет и трифазна статорска намотка и нема
механички комутатор и четки преку кои конвенционалниот еднонасочен мотор ја
префрла струјата меѓу намотките додека роторот се врти~\cite{onsemi_bldc,nxp_backemf}.
Функцијата на комутација што кај четкастиот мотор ја вршат четките механички затоа
мора да се изврши електронски. Инверторот е токму колото што го прави тоа: го
претвора еднонасочниот напон на сабирницата со фиксен поларитет (овде обезбеден од
батерија) во низа контролирани наизменични фазни струи потребни за да се одржи
статорското магнетно поле во вртење пред роторот и со тоа да се произведува
континуиран вртежен момент~\cite{onsemi_bldc,nxp_backemf}.

Во трапезоидното (шестчекорно, \ang{120}) управување --- најчестата BLDC шема ---
шесте прекинувачи на инверторот заземаат шест комбинации на вклучено/исклучено,
произведувајќи шест дискретни ориентации на статорското поле по електрична
револуција; прекинувачите се чекорат низ оваа низа синхроно со роторот, така што
вртежниот момент е секогаш позитивен~\cite{onsemi_bldc}. Положбата на роторот ---
потребна за да се знае \emph{кога} да се чекори --- се добива или од Hall сензори
поставени на \ang{120}, или, безсензорски, од back-EMF (контраелектромоторната
сила) на моментно неенергизираната фаза~\cite{onsemi_bldc,nxp_backemf}. Во секој
момент една фаза се управува позитивно, една негативно, а трета се остава отворена,
совпаѓајќи ја правоаголната струја со трапезоидната back-EMF на моторот за ефикасно
производство на момент~\cite{onsemi_bldc}. PWM применет на спроводните прекинувачи
го поставува ефективниот напон, а со тоа и работната точка по брзина и момент~\cite{nxp_backemf}.
Без инвертор нема начин да се комутира безчеткеста машина --- токму затоа
енергетската електроника, а не четките, го дефинира погонот.

\subsection{Како работи инверторското коло}
Енергетскиот степен е трифазен мост: три идентични „гранки“ (полумостови), по една
за секоја фаза на моторот, секоја составена од high-side и low-side прекинувач
поврзани во серија преку еднонасочната сабирница, со фазата на моторот приклучена
на нивната средна точка. Во овој дизајн секој прекинувач е составен од два паралелно
поврзани IPT015N10N5 транзистори, па една гранка содржи четири MOSFET-а, а целиот
мост дванаесет.

Во рамките на една гранка двата прекинувачи работат комплементарно --- никогаш и
двата истовремено вклучени, што би ја кратко споило сабирницата (shoot-through).
Вклучувањето на high-side прекинувачот ја поврзува фазата со позитивната сабирница;
вклучувањето на low-side прекинувачот ја поврзува со масата; оставањето на двата
исклучени ѝ дозволува на фазата да лебди (што се користи за неенергизираната фаза и
за следење на back-EMF)~\cite{onsemi_bldc,nxp_backemf}. Чекорењето на трите гранки
низ шестчекорната шема ја циркулира струјата низ парови фази по редоследот што
произведува ротирачко статорско поле~\cite{onsemi_bldc}. За време на dead-time-от
меѓу состојбите на прекинувачите, телесните диоди (body diode) на MOSFET-ите ---
заедно со индуктивноста на моторот --- ја обезбедуваат freewheeling патеката за
индуктивната фазна струја; телесната диода на IPT015N10N5 е карактеризирана токму за
оваа намена ($V_{SD}\approx\SI{0.85}{\volt}$, $Q_{rr}\approx\SI{316}{\nano\coulomb}$)~\cite{ipt015ds}.
Отпорничкиот карактер на спроводната состојба на MOSFET-от ($R_{DS(on)}$ од само
\SI{1.5}{\milli\ohm}, дополнително преполовена со паралелното поврзување) ги држи
спроводните загуби ниски при струите од редот на \SI{100}{\ampere}~\cite{ipt015ds,balogh}.

\subsection{Намена на секоја компонента}
\begin{itemize}[leftmargin=1.3em]
  \item \textbf{Енергетски MOSFET-и (IPT015N10N5, $\times 2$ по позиција):}
        контролираните прекинувачи што ја поврзуваат секоја фаза со сабирниците.
        Двата паралелно поврзани уреди ја преполовуваат ефективната отпорност во
        спроводна состојба и ги делат загубата и топлината; нивниот позитивен
        температурен коефициент на $R_{DS(on)}$ го прави делењето
        самостабилизирачко~\cite{balogh}.
  \item \textbf{Gate driver (UCC27211A-Q1):} ја претвора PWM логиката на
        контролерот во високострујното \SI{12}{\volt} управување на гејтот потребно
        за брзо прекинување на MOSFET-ите, го поместува нивото на high-side
        сигналот до лебдечкиот јазол на прекинување и спроведува UVLO
        заштита~\cite{ucc27211}. Драјверите постојат токму затоа што логичкиот излез
        (често \SI{3.3}{\volt}, $\le\SI{1}{\ampere}$) не може ниту целосно да
        спроведе енергетски MOSFET, ниту да ја испорача врвната струја за брзо
        прекинување, и затоа што гејтот на high-side уредот мора да биде референциран
        кон јазол што се ниша до напонот на сабирницата~\cite{ucc27211,balogh}.
  \item \textbf{Bootstrap кондензатор:} лебдечкиот резервоар на енергија што го
        напојува high-side драјверот додека јазолот на прекинување е висок
        (в.\ Дел~\ref{sec:boottheory}).
  \item \textbf{Bootstrap диода (внатрешна):} еднонасочен вентил што го надополнува
        bootstrap кондензаторот од $V_{DD}$ кога јазолот на прекинување е низок и го
        блокира напонот на сабирницата кога е висок~\cite{ucc27211,infineon_monolithic}.
  \item \textbf{Отпорници на гејтот:} ја поставуваат и ограничуваат струјата на
        гејтот, го наштелуваат $\mathrm{d}v/\mathrm{d}t$ заради EMI и ги пригушуваат
        и балансираат паралелните гејтови~\cite{balogh}.
  \item \textbf{Pull-down отпорници гејт--извор:} ги држат MOSFET-ите исклучени при
        вклучување на напојувањето и додаваат имунитет на вклучување од
        $\mathrm{d}v/\mathrm{d}t$~\cite{balogh}.
  \item \textbf{Кондензатори за декуплирање:} ги напојуваат брзите струјни импулси
        на гејтот локално, така што \SI{12}{\volt} сабирницата и преднапонот на
        драјверот не се рушат при прекинување~\cite{ucc27211}.
  \item \textbf{Капацитивност на DC-link:} ја вкрутува сабирницата против големата
        бранувачка струја на полумостот и го ограничува надскокот при комутација во
        рамки на оценката од \SI{100}{\volt} на уредот.
\end{itemize}

\subsection{Bootstrap колото: работа, потреба и улогата на Miller plateau}
\label{sec:boottheory}
\paragraph{Зошто полумост само со N-канални транзистори бара bootstrap напојување.}
N-каналните MOSFET-и се претпочитаат за полумостот бидејќи, за дадена површина на
чипот и напон, имаат многу пониска $R_{DS(on)}$ од P-каналните уреди. Но N-канален
MOSFET се вклучува само кога неговиот \textbf{гејт е неколку волти над неговиот
извор}. Кај low-side уредот изворот е на маса, па \SI{12}{\volt} управување
референцирано кон маса работи директно. Кај high-side уредот изворот е јазолот на
прекинување, кој се крева до полниот напон на сабирницата кога тој уред е вклучен;
за да остане вклучен, гејтот мора да се управува $\approx\SI{12}{\volt}$
\emph{над сабирницата} --- напон што не постои никаде во систем со една
сабирница~\cite{ucc27211,balogh}. Bootstrap колото го создава ова лебдечко
high-side напојување евтино, користејќи само диода и кондензатор~\cite{infineon_monolithic}.

\paragraph{Како работи.} Кога low-side прекинувачот е вклучен, јазолот на
прекинување (и пинот HS) се влече кон маса; bootstrap кондензаторот, поврзан меѓу HB
и HS, се полни од \SI{12}{\volt} сабирницата $V_{DD}$ низ bootstrap диодата до
$\approx V_{DD}-V_F\approx\SI{11}{\volt}$~\cite{ucc27211,infineon_monolithic}. Кога
потоа high-side-от се вклучува, јазолот на прекинување се крева до напонот на
сабирницата; bootstrap диодата станува обратно поларизирана (го блокира напонот на
сабирницата), а наполнетиот кондензатор „се вози нагоре“ заедно со јазолот на
прекинување, држејќи го пинот HB $\approx\SI{11}{\volt}$ над HS и така
обезбедувајќи го лебдечкиот преднапон што го држи high-side MOSFET-от целосно
спроводен~\cite{ucc27211,infineon_monolithic}. Кондензаторот повторно се дополнува
во секој циклус кога low-side-от следно се вклучува --- поради што е потребно
минимално време на спроводност на low-side-от (освежување) и поради што
Дел~\ref{sec:cboot} мораше да ги провери и опаѓањето и полнењето~\cite{infineon_monolithic,infineon_bootoverview}.

\paragraph{Каде влегува Miller plateau.} За време на секое вклучување напонот на
гејтот расте, паузира на рамното „\eng{Miller plateau}“ ($\approx\SI{4.4}{\volt}$
овде~\cite{ipt015ds}), потоа повторно расте. Паузата настанува бидејќи, штом каналот
ја носи полната струја на оптоварување, целата струја за управување на гејтот се
пренасочува во празнење на капацитивноста гејт--одвод $C_{GD}$, така што напонот на
одводот може да се ниша; во текот на овој интервал $V_{GS}$ е стегнат и се движи
само напонот на одводот~\cite{balogh}. Интервалот на платото (заедно со интервалот
од прагот до платото) е токму прозорецот во кој уредот истовремено издржува висок
напон и носи висока струја, па тоа е местото каде што се генерира загубата при
прекинување~\cite{balogh}. Miller полнежот $Q_{gd}$ што мора да се протурка низ
платото е удвоен со паралелното поврзување (овде $\approx\SI{102}{\nano\coulomb}$ во
најлош случај), а \emph{брзината} на напонската транзиција ја поставува тоа колку
струја драјверот може да втисне во гејтот на нивото на платото --- а не неговата
врвна струја при полн $V_{DD}$~\cite{balogh}.

\paragraph{Зошто се потребни големи изворни и понорни струи.} Анализата на Balogh го
истакнува тоа експлицитно: најважната карактеристика на еден gate driver е неговата
способност за извор и понор \emph{околу напонот на Miller plateau-то}
($\sim\SI{5}{\volt}$), бидејќи таа струја одредува колку брзо се одвива напонската
транзиција на одводот и оттаму загубата при прекинување~\cite{balogh}. Голема
\textbf{изворна} струја (\SI{3.7}{\ampere}) брзо го управува гејтот низ платото при
вклучување, рушејќи го $V_{DS}$ за $\approx\SI{28}{\nano\second}$
(Дел~\ref{sec:gateR}) и минимизирајќи ја загубата при вклучување~\cite{ucc27211,balogh}.
Голема \textbf{понорна} струја (\SI{4.5}{\ampere}) го прави истото при исклучување
и, што е клучно, го држи гејтот на \emph{исклучениот} уред цврсто на маса додека
спротивниот прекинувач го ниша споделениот јазол: силниот понор ја надвладува Miller
струјата инјектирана низ $C_{GD}$, спречувајќи $\mathrm{d}v/\mathrm{d}t$-индуцирано
(Miller) вклучување и резултирачкиот shoot-through~\cite{ucc27211,balogh}. Со два
паралелни MOSFET-и со висок $Q_g$ по позиција, токму оваа висока струјна способност
прави еден канал на драјверот да биде доволен --- и затоа TI ја пласира струјата на
уредот специфично за „управување на големи енергетски MOSFET-и \ldots{} низ Miller
plateau-то“~\cite{ucc27211}.

\section{Спецификација на дизајнот и работен опсег}\label{sec:env}
Зададените параметри на дизајнот се сумирани во Табела~\ref{tab:spec}. Хардверот
за енергетскиот степен мора да биде димензиониран за најголемата струја што
инверторот може да ја носи; истата струја ги одредува вредноста на shunt-от и
засилувањето на струјниот засилувач. Бидејќи моторот се напојува од батерија со
променлив напон, струите за најлош случај мора да се проценат при \emph{најнискиот}
напон на пакетот, зашто за дадена моќност струјата расте кога напонот опаѓа
($I=P/V$).

\begin{table}[t]
\centering
\caption{Зададени параметри на дизајнот.}
\label{tab:spec}
\small
\begin{tabularx}{\columnwidth}{@{}>{\RaggedRight\arraybackslash}p{3.4cm} >{\RaggedRight\arraybackslash}p{3.6cm} L@{}}
\toprule
\textbf{Параметар} & \textbf{Вредност} & \textbf{Образложение / извор} \\
\midrule
Извор на напојување & Батерии (\SI{24}{\volt} или \SI{48}{\volt} пакет) &
Напонот варира со состојбата на полнење на батеријата \\
Напон на DC-сабирницата (DC-link) & \SI{24}{\volt} пакет: \mbox{20--\SI{30}{\volt}}
\newline \SI{48}{\volt} пакет: \mbox{40--\SI{60}{\volt}} & $V_{bus,max}=\SI{60}{\volt}$
(полна \SI{48}{\volt} батерија), $V_{bus,min}=\SI{20}{\volt}$ (празна \SI{24}{\volt}
батерија) \\
Излезна моќност & $\approx\SI{3}{\kilo\watt}$ (\SI{24}{\volt} пакет) /
$\approx\SI{5}{\kilo\watt}$ (\SI{48}{\volt} пакет) & При номинален напон \\
Склопна фреквенција, $f_{SW}$ & \SI{40}{\kilo\hertz} & Над границата на човечкиот
слух ($\sim\SI{20}{\kilo\hertz}$), па PWM тонот е нечуен \\
Прекинувачки уред & IPT015N10N5, $\times 2$ паралелно по позиција &
\SI{100}{\volt}, \SI{1.5}{\milli\ohm} OptiMOS\textsuperscript{\texttrademark}~5,
оптимизиран за апликации со голема струја~\cite{ipt015ds,ipt015part} \\
Gate driver & UCC27211A-Q1, $\times 1$ по фаза & \SI{120}{\volt} boot,
\SI{3.7}{\ampere}/\SI{4.5}{\ampere} полумостен драјвер со вградена boot
диода~\cite{ucc27211,ucc27211folder} \\
Напон за управување на гејтот, $V_{DD}$ & \SI{12}{\volt} & Во работниот опсег
\mbox{8--\SI{17}{\volt}} на уредот~\cite{ucc27211} \\
\bottomrule
\end{tabularx}
\end{table}

\subsection{Струја на сабирницата за најлош случај (термичко димензионирање)}
Конзервативна граница се добива ако се претпостави дека конверторот сè уште бара
номинална моќност и при празна батерија:
\begin{itemize}[leftmargin=1.4em]
  \item пакет \SI{24}{\volt} (празен, \SI{20}{\volt}), \SI{3}{\kilo\watt}
        $\rightarrow I_{bus,max}=3000/20=\SI{150}{\ampere}$;
  \item пакет \SI{48}{\volt} (празен, \SI{40}{\volt}), \SI{5}{\kilo\watt}
        $\rightarrow I_{bus,max}=5000/40=\SI{125}{\ampere}$.
\end{itemize}
Меродавната струја за термичко димензионирање затоа е $\approx\SI{150}{\ampere}$
(\SI{24}{\volt} пакет при празна батерија); за споредба, при номинален напон струите
се $5000/48\approx\SI{104}{\ampere}$ и $3000/24\approx\SI{125}{\ampere}$. Бидејќи во
секоја позиција се поврзани паралелно два уреди, секој MOSFET носи приближно
половина од оваа струја, т.е.\ до $\approx\SI{75}{\ampere}$ во најлош случај, што е
далеку под континуираната оценка од $2\times\SI{250}{\ampere}$ ($T_C=\SI{100}{\celsius}$)
на паралелниот пар~\cite{ipt015ds}. Спроводната загуба по уред во овој најлош
случај е $P=I^2R_{DS(on)}=75^2\times\SI{1.5}{\milli\ohm}\approx\SI{8.4}{\watt}$
(моментно, при спроводност), што бара соодветно ладење.

\subsection{Фазна струја за најлош случај (мерење на струја)}
Електричната влезна моќност $P$ е еднаква на излезната AC моќност на инверторот
поделена со ефикасноста на конверторот $\eta$; земањето $\eta\!\to\!1$ е
конзервативната претпоставка (што ја максимизира струјата). За синусоидно модулиран
трифазен инвертор со просторно-векторска PWM (SVPWM), максималната амплитуда на
основната хармоника на фазниот напон е $\hat V_{\mathrm{ph}}=\Vdcbus/\sqrt{3}$, па
максималниот RMS фазен напон е
\begin{equation}
V_{\mathrm{ph,rms}}=\frac{\Vdcbus}{\sqrt{6}} .
\label{eq:vph}
\end{equation}
Со три фази и фактор на моќност на поместување $\mathrm{PF}=\cos\varphi$, AC моќноста
е $P=3\,V_{\mathrm{ph,rms}}\,\Iphrms\,\mathrm{PF}$. Со замена на~\eqref{eq:vph} и
решавање по RMS фазната струја,
\begin{equation}
\boxed{\;\Iphrms=\frac{\sqrt{6}}{3}\,\frac{P}{\Vdcbus\,\mathrm{PF}}
        \approx 0.816\,\frac{P}{\Vdcbus\,\mathrm{PF}}\;}
\label{eq:iph}
\end{equation}
која е најголема при минимален напон на сабирницата и минимален фактор на моќност.
Равенката~\eqref{eq:iph} ја дава струјата во номиналната работна точка; треба да се
забележи дека мотор при ниска брзина и висок вртежен момент може да бара споредлива
фазна струја при \emph{пониска} моќност на сабирницата, па~\eqref{eq:iph} се користи
како термичка/континуирана проектна точка, а дополнителна резерва се остава за
преодни и stall услови во опсегот на засилувачот (Дел~\ref{sec:csa}).

\begin{table}[t]
\centering
\caption{Работни точки и резултантна фазна струја од~\eqref{eq:iph} ($\mathrm{PF}=0.9$).}
\label{tab:wcphase}
\small
\begin{tabular}{l c c c c}
\toprule
Напојување / мотор & \(\Vdcbus\) (V) & \(P\) (kW) & \(\Iphrms\) (A,rms) & \(\hat I_{\mathrm{ph}}\) (A,врв)\\
\midrule
\SI{24}{\volt} бат., \SI{3}{\kilo\watt} & 20 (мин) & 3 & \textbf{136.1} & \textbf{192.5}\\
\SI{24}{\volt} бат., \SI{3}{\kilo\watt} & 30 (макс) & 3 & 90.7 & 128.3\\
\SI{48}{\volt} бат., \SI{5}{\kilo\watt} & 40 (мин) & 5 & 113.4 & 160.4\\
\SI{48}{\volt} бат., \SI{5}{\kilo\watt} & 60 (макс) & 5 & 75.6 & 106.9\\
\bottomrule
\end{tabular}
\end{table}

Ограничувачката работна точка е \SI{3}{\kilo\watt} при \SI{20}{\volt}:
$\Iphrms^{\max}\approx\SI{136}{\ampere}$ RMS, $\hat I_{\mathrm{ph}}^{\max}\approx\SI{193}{\ampere}$
врв (Табела~\ref{tab:wcphase}). Ова е под номиналната граница од \SI{300}{\ampere}
за пакувањето на IPT015N10N5~\cite{ipt015part}, што потврдува дека енергетскиот
степен е во согласност и со фазната струјна цел. Сите вредности за мерење на струја
во Дел~\ref{sec:current} се однесуваат на оваа работна точка.

\section{Избор и компатибилност на енергетските компоненти}\label{sec:power}

\subsection{Infineon IPT015N10N5 --- клучни параметри}
IPT015N10N5 е N-канален OptiMOS\textsuperscript{\texttrademark}~5 енергетски
транзистор во TO-Leadless (PG-HSOF-8, „TOLL“) куќиште, кое Infineon експлицитно го
опишува како идеално за прекинување со висока фреквенција и наменето за погони на
лесни електрични возила, виљушкари и електрични алати~\cite{ipt015ds,ipt015part}.
Параметрите релевантни за овој дизајн се сумирани во Табела~\ref{tab:ipt}; сите
вредности се од каталошкиот лист рев.\,2.4, освен ако не е поинаку наведено~\cite{ipt015ds}.

\begin{table}[t]
\centering
\caption{Клучни параметри на IPT015N10N5~\cite{ipt015ds}.}
\label{tab:ipt}
\small
\begin{tabularx}{\columnwidth}{@{}l L L@{}}
\toprule
\textbf{Симбол} & \textbf{Параметар} & \textbf{Вредност / забелешка} \\
\midrule
$V_{DS}$ & Напон одвод--извор & \SI{100}{\volt} (апсолутен максимум) \\
$R_{DS(on)}$ & Отпорност во спроводна состојба ($V_{GS}=\SI{10}{\volt}$) &
\SI{1.3}{\milli\ohm} тип.\ / \SI{1.5}{\milli\ohm} макс. \\
$I_D$ & Континуирана струја на одвод & \SI{353}{\ampere} ($T_C=\SI{25}{\celsius}$);
\SI{250}{\ampere} ($T_C=\SI{100}{\celsius}$) \\
$V_{GS(th)}$ & Прагов напон на гејтот & 2.2 / 3.0 / \SI{3.8}{\volt} (мин./тип./макс.) \\
$V_{plateau}$ & Напон на Miller plateau & \SI{4.4}{\volt} тип.\ (при $I_D=\SI{100}{\ampere}$) \\
$C_{iss}$ & Влезна капацитивност & \SI{12}{\nano\farad} тип.\ / \SI{16}{\nano\farad} макс. \\
$C_{oss}$ & Излезна капацитивност & \SI{1.8}{\nano\farad} тип. \\
$C_{rss}$ & Преносна (reverse-transfer) капацитивност & \SI{80}{\pico\farad} тип.\ / \SI{140}{\pico\farad} макс. \\
$R_{G}$ & Внатрешна отпорност на гејтот & \SI{1.4}{\ohm} тип.\ / \SI{2.1}{\ohm} макс. \\
$Q_{g}$ & Вкупен полнеж на гејтот (\mbox{0--\SI{10}{\volt}}) & \SI{169}{\nano\coulomb} тип.\ / \SI{211}{\nano\coulomb} макс. \\
$Q_{gs}$ & Полнеж гејт--извор & \SI{53}{\nano\coulomb} \\
$Q_{gd}$ & Полнеж гејт--одвод (Miller) & \SI{34}{\nano\coulomb} тип.\ / \SI{51}{\nano\coulomb} макс. \\
$g_{fs}$ & Трансконтуктанса & \SI{140}{\siemens} мин.\ / \SI{280}{\siemens} тип. \\
$V_{SD}$ & Напон на телесната диода (body diode) & \SI{0.85}{\volt} тип.\ / \SI{1.2}{\volt} макс. \\
$V_{GS,max}$ & Граничен напон гејт--извор & $\pm\SI{20}{\volt}$ \\
\bottomrule
\end{tabularx}
\end{table}

Особина што директно го оправдува паралелното поврзување е позитивниот температурен
коефициент на $R_{DS(on)}$ (околу $+0.7$ до $+1\,\%/\si{\celsius}$ кај енергетските
MOSFET-и): уред што спроведува поголема струја повеќе се загрева, неговата отпорност
расте и струјата се прелева кон соседниот уред, така што паралелните MOSFET-и
самостојно ги изедначуваат своите струи~\cite{balogh}. Ова е стандардната причина
зошто MOSFET-ите --- за разлика од биполарните транзистори --- се добро прилагодени
за паралелна работа~\cite{balogh}.

\subsection{Texas Instruments UCC27211A-Q1 --- клучни параметри}
UCC27211A-Q1 е AEC-Q100 квалификуван \SI{120}{\volt} high-side/low-side полумостен
gate driver, дизајниран да управува два N-канални MOSFET-и со независни влезови;
неговиот лебдечки high-side степен може да работи до \SI{120}{\volt}, а
\SI{120}{\volt} bootstrap диода е интегрирана на чипот, така што надворешна boot
диода не е задолжителна~\cite{ucc27211,ucc27211folder}. Релевантните параметри, сите
од каталошкиот лист SLUSCG0C, се дадени во Табела~\ref{tab:ucc}~\cite{ucc27211}.

\begin{table}[t]
\centering
\caption{Клучни параметри на UCC27211A-Q1~\cite{ucc27211}.}
\label{tab:ucc}
\small
\begin{tabularx}{\columnwidth}{@{}l L L@{}}
\toprule
\textbf{Симбол} & \textbf{Параметар} & \textbf{Вредност} \\
\midrule
$V_{DD}$ & Напојување (работно / апс.\ макс.) & \mbox{8--\SI{17}{\volt}} / \SI{20}{\volt} \\
$I_{O,source}$ / $I_{O,sink}$ & Врвна излезна струја & \SI{3.7}{\ampere} / \SI{4.5}{\ampere} \\
$V_{DDR}$ & UVLO праг на $V_{DD}$ (растечки) & \SI{7.0}{\volt} (хистереза \SI{0.5}{\volt}) \\
$V_{HBR}$ & UVLO праг на high-side (HB--HS, растечки) & \SI{6.7}{\volt} (хистереза \SI{1.1}{\volt}) \\
$V_{HB}$ & Макс.\ напон на boot напојувањето & \SI{120}{\volt} (HB--HS макс.\ \SI{20}{\volt}) \\
$V_{HS}$ & Опсег на јазолот на прекинување & $-(24-V_{DD})$\,V ($-\SI{12}{\volt}$ при \SI{12}{\volt}), до \SI{105}{\volt} (препорачано) \\
$V_{F}$ & Пад на напон на boot диодата & \SI{0.85}{\volt} (@\SI{100}{\micro\ampere}) / \SI{1.05}{\volt} (@\SI{100}{\milli\ampere}) макс. \\
$R_{D}$ & Динамичка отпорност на boot диодата & \SI{0.55}{\ohm} тип. \\
$t_{PD}$ & Доцнење на простирање (propagation delay) & \SI{20}{\nano\second} тип. \\
$t_{R}$/$t_{F}$ & Време на пораст/опаѓање (\SI{1}{\nano\farad}) & \SI{7.2}{\nano\second} / \SI{5.5}{\nano\second} \\
$I_{HB}$ & Струја на мирување на boot & \SI{0.12}{\milli\ampere} макс. \\
$SR_{HS}$ & Макс.\ slew rate на HS & \SI{50}{\volt\per\nano\second} \\
$C_{HB}$ & Препорачан опсег на boot кондензатор & \mbox{0.022--\SI{0.1}{\micro\farad}} (зависно од полнежот на гејтот) \\
$C_{VDD}$ & Препорачано декуплирање на $V_{DD}$ & \mbox{0.22--\SI{4.7}{\micro\farad}} \\
\bottomrule
\end{tabularx}
\end{table}

\subsection{Внатрешна архитектура на UCC27211A-Q1}
TI го дели уредот на пет функционални блока: влезни степени, UVLO заштита,
поместување на ниво (level shift), boot диода и излезни драјверски степени~\cite{ucc27211}.
\begin{itemize}[leftmargin=1.3em]
  \item \textbf{Влезни степени.} Независни HI и LI влезови со $\approx\SI{68}{\kilo\ohm}$
        pull-down, $\approx\SI{4}{\pico\farad}$ влезна капацитивност и
        TTL-компатибилни прагови (растечки \SI{2.3}{\volt}, опаѓачки \SI{1.6}{\volt})
        со широка хистереза за имунитет на шум; влезовите поднесуваат
        \mbox{$-10$ до $+\SI{20}{\volt}$} и се независни од $V_{DD}$, па прифаќаат и
        логика од микроконтролер и сигнали од аналоген контролер~\cite{ucc27211}.
  \item \textbf{UVLO заштита.} Посебни монитори на $V_{DD}$ и на лебдечкото
        $V_{HB}-V_{HS}$ напојување ги принудуваат излезите ниски додека преднапонот
        не е валиден. Прагот на $V_{DD}$ е \SI{7}{\volt} (хистереза \SI{0.5}{\volt}),
        а симетричниот high-side праг е \SI{6.7}{\volt} (хистереза \SI{1.1}{\volt}),
        така што high-side-от и low-side-от се исклучуваат заедно~\cite{ucc27211}.
  \item \textbf{Поместување на ниво (level shift).} Ја преведува командата на
        high-side влезот, референцирана кон маса, нагоре до лебдечкиот high-side
        излезен степен референциран кон HS, со тесно ($\approx\SI{4}{\nano\second}$
        типично) усогласување на доцнењето меѓу каналите, за да се избегне временска
        нерамнотежа во полумостот~\cite{ucc27211}.
  \item \textbf{Boot диода.} \SI{120}{\volt} bootstrap диодата е интегрирана (анода
        $V_{DD}$, катода HB), елиминирајќи надворешна диода и освежувајќи го
        $C_{BOOT}$ во секој циклус кога HS оди ниско~\cite{ucc27211}.
  \item \textbf{Излезни степени.} Low-side степенот е референциран
        $V_{DD}\rightarrow V_{SS}$, а high-side степенот $V_{HB}\rightarrow V_{HS}$;
        двата обезбедуваат \SI{3.7}{\ampere} извор / \SI{4.5}{\ampere} понор со брзи
        (\SI{7.2}{\nano\second} пораст / \SI{5.5}{\nano\second} опаѓање во
        \SI{1}{\nano\farad}) излези со ниска импеданса и доцнење на простирање од
        \SI{20}{\nano\second}~\cite{ucc27211}.
\end{itemize}

\paragraph{Проектни равенки од каталошкиот лист.} Дисипацијата на драјверот се дели
на еднонасочен и прекинувачки дел~\cite{ucc27211}:
\begin{equation}
P_{DISS}=P_{DC}+P_{SW},\qquad P_{DC}=I_Q\times V_{DD},
\end{equation}
со $I_Q<\SI{0.17}{\milli\ampere}$, па $P_{DC}$ е занемарлив~\cite{ucc27211}.
Прекинувачкиот дел се изведува од енергијата за полнење на гејтот, претставена како
еквивалентен кондензатор~\cite{ucc27211}:
\begin{equation}
E_G=\tfrac{1}{2}\,C_{LOAD}\,V_{DD}^{2},\qquad
P_G=C_{LOAD}\,V_{DD}^{2}\,f_{SW}=Q_G\,V_{DD}\,f_{SW},
\end{equation}
каде што гејтот на MOSFET-от е претставен преку неговиот полнеж низ
$Q_G=C_{LOAD}\,V_{DD}$~\cite{ucc27211}. Половина од $P_G$ се дисипира при вклучување,
а половина при исклучување; со надворешни отпорници на гејтот оваа загуба се дели
меѓу драјверот и отпорниците~\cite{ucc27211}. Истиот израз се појавува кај Balogh
како загуба за управување на гејтот, $P_{GATE}=V_{DRV}\,Q_G\,f_{DRV}$, чиј член
$Q_G f_{DRV}$ е просечната струја на преднапон на гејтот~\cite{balogh}. За овој
дизајн, моќноста за управување на гејтот по драјвер (двата прекинувачи на една
гранка, \SI{12}{\volt}, \SI{40}{\kilo\hertz}, $\approx\SI{465}{\nano\coulomb}$ секој)
е приближно
\begin{equation}
P_G\approx 2\times Q_{g,\text{поз}}\,V_{DD}\,f_{SW}
        =2\times\SI{465}{\nano\coulomb}\times\SI{12}{\volt}\times\SI{40}{\kilo\hertz}
        \approx\SI{0.45}{\watt}\ \text{по драјвер,}
\end{equation}
поделена меѓу драјверот и осумте отпорници на гејтот на гранката --- удобно во рамки
на дисипацијата на SO-PowerPAD куќиштето кога е изведено според упатствата на TI
(драјвер близу MOSFET-ите, декуплирање директно кај пиновите, тесни кола на
гејтот)~\cite{ucc27211}. Барањето за врвна струја е сопствениот проектен критериум на
каталошкиот лист: драјверот мора да го испорача Miller полнежот во рамки на целното
време на транзиција, $I_{PEAK}=Q_{GD}/\mathrm{d}t$, и способноста од \SI{3.7}{\ampere}
е проверена наспроти ова во Дел~\ref{sec:gateR}~\cite{ucc27211}.

\subsection{Заклучок за компатибилноста}
Табела~\ref{tab:compat} ги сумира проверките на компатибилноста за целиот опсег на
напони на батериите.

\begin{table}[t]
\centering
\caption{Проверка на компатибилноста на компонентите.}
\label{tab:compat}
\small
\begin{tabularx}{\columnwidth}{@{}L L L >{\RaggedRight\arraybackslash}p{2.0cm}@{}}
\toprule
\textbf{Проверка} & \textbf{Барање} & \textbf{Овој дизајн} & \textbf{Исход} \\
\midrule
Напон на блокирање & $V_{DS}\gg V_{bus,max}$ вкл.\ ringing & \SI{100}{\volt}
наспроти \SI{60}{\volt}; $\approx 1.67\times$ DC маргина & \textbf{Поминува} ---
$\approx\SI{40}{\volt}$ резерва за надскок при исклучување~\cite{ipt015ds} \\
Напон на high-side напојувањето & $V_{HB}>V_{bus,max}$ & \SI{120}{\volt} boot,
\SI{105}{\volt} препорачано HS, наспроти \SI{60}{\volt} & \textbf{Поминува}~\cite{ucc27211} \\
Ниво на управување на гејтот & $V_{GS,max}>V_{DD}\ge$ полна спроводливост &
$\pm\SI{20}{\volt}$ граница, управуван на \SI{12}{\volt} & \textbf{Поминува}~\cite{ipt015ds,ucc27211} \\
UVLO & $V_{DD}>$ UVLO & \SI{12}{\volt} наспроти \SI{7}{\volt} ($V_{DD}$),
\SI{6.7}{\volt} (HB) & \textbf{Поминува}~\cite{ucc27211} \\
Негативно ringing на јазолот & HS поднесува потсек & $-\SI{12}{\volt}$ преодна
способност при $V_{DD}=\SI{12}{\volt}$ & \textbf{Поминува}~\cite{ucc27211} \\
Струен капацитет & $I_D\gg$ струја на гранката & $2\times\SI{250}{\ampere}$
(\SI{100}{\celsius}) наспроти $\approx\SI{75}{\ampere}$ (најлош случај при
$V_{bus,min}=\SI{20}{\volt}$) & \textbf{Поминува}~\cite{ipt015ds} \\
Сила на управување наспроти полнеж & Доволна струја низ Miller plateau за
\emph{паралелниот} полнеж & в.\ Дел~\ref{sec:gateR} & \textbf{Поминува} \\
\bottomrule
\end{tabularx}
\end{table}

Спрегата е електрично исправна за целиот опсег на напони на батериите. При
највисокиот напон на сабирницата (\SI{60}{\volt} при полна \SI{48}{\volt} батерија)
\SI{100}{\volt} класата дава DC маргина од $\approx 1.67\times$ и околу
\SI{40}{\volt} резерва за надскокот на $V_{DS}$ при исклучување предизвикан од
$L\cdot\mathrm{d}i/\mathrm{d}t$; ова е поудобно отколку што изгледа, но го прави
управувањето на брзината на исклучување (преку асиметричната мрежа на отпорници во
Дел~\ref{sec:asym}) уште поважно отколку на фиксна \SI{48}{\volt} сабирница.
Единствената точка што бара квантитативна проверка --- бидејќи дизајнот паралелно
поврза два MOSFET-и со висок $Q_g$ по прекинувач --- е дали еден канал на
UCC27211A-Q1 може доволно брзо да го наполни и испразни \emph{удвоениот} полнеж на
гејтот. Ова е разгледано во Дел~\ref{sec:gateR} и во димензионирањето на bootstrap-от
во Дел~\ref{sec:cboot} и е причината зошто стандардниот \SI{0.1}{\micro\farad}
bootstrap кондензатор препорачан во каталошкиот лист мора да се преиспита за оваа
оптоварувачка состојба.

\section{Дизајн на колото за управување на гејтот}\label{sec:gatedrive}
Низ целиот овој дел, полнежот на гејтот за една \emph{прекинувачка позиција} е двојно
поголем од оној на еден уред, бидејќи се поврзани паралелно два MOSFET-и:
\begin{equation}
Q_{g,\text{поз}}=2\times Q_{g}=2\times\SI{211}{\nano\coulomb}=\SI{422}{\nano\coulomb}
\quad\text{(најлош случај, до \SI{10}{\volt})}~\cite{ipt015ds}.
\end{equation}
Бидејќи гејтот се управува до \SI{12}{\volt} наместо до \SI{10}{\volt} (при кои е
специфициран $Q_g$), и бидејќи кривата на полнежот е речиси рамна над
$\sim\SI{7}{\volt}$ (па дополнителниот полнеж од 10 до \SI{12}{\volt} е претежно
линеарниот член $C_{GS}$), се додава резерва од $\approx 10\,\%$. Дополнително,
$Q_g$ е специфициран при напон одвод--извор од \SI{50}{\volt}~\cite{ipt015ds}; при
најлошиот напон на сабирницата од \SI{60}{\volt}, Miller полнежот $Q_{gd}$ (кој
зависи од нишањето на напонот на одводот) е маргинално поголем за неколку nC по уред
(бидејќи $C_{rss}$ нагло опаѓа со напонот), но и тоа е опфатено со истата резерва.
Резултирачката проектна вредност е
\begin{equation}
Q_{g,\text{поз,проект}}\approx\SI{465}{\nano\coulomb}\ \text{по прекинувачка позиција.}
\end{equation}

\subsection{Bootstrap кондензатор}\label{sec:cboot}
\paragraph{Метод.} Bootstrap кондензаторот $C_{BOOT}$ го обезбедува целиот полнеж
потребен за вклучување на high-side уредот и за напојување на струјата на мирување на
high-side степенот во текот на спроводноста на high-side-от, додека опаѓа за не
повеќе од дозволената вредност $\Delta V_{BS}$~\cite{infineon_monolithic,infineon_bootcalc}.
Управувачките равенки се
\begin{align}
\Delta V_{BS} &\le V_{DD}-V_{F}-V_{GS,\min}-V_{DS,on}, \label{eq:dvbs}\\
Q_{TOTAL} &=Q_{g,\text{поз}}+\frac{I_{QBS}+I_{LK}}{f_{SW}}+Q_{LS},\qquad
C_{BOOT}\ge\frac{Q_{TOTAL}}{\Delta V_{BS}}, \label{eq:cboot}
\end{align}
со дополнителна резерва од $\sim 20\,\%$ препорачана за толеранции и за опаѓањето на
капацитивноста кога диодата е интегрирана~\cite{infineon_maxboot}.

\paragraph{Бројки.}
\begin{itemize}[leftmargin=1.4em]
  \item Пад на boot диодата при релевантната струја на полнење:
        $V_F\approx\SI{0.9}{\volt}$~\cite{ucc27211}.
  \item Се бара гејтот на high-side-от да остане целосно спроводен,
        $V_{GS,\min}\approx\SI{10}{\volt}$; low-side $V_{DS,on}\approx\SI{0.1}{\volt}$.
        Тогаш \emph{дозволеното} опаѓање според~\eqref{eq:dvbs} е
        $\Delta V_{BS}\le 12-0.9-10-0.1=\SI{1.0}{\volt}$~\cite{infineon_monolithic}.
        Се избира конзервативна проектна цел од \textbf{$\Delta V_{BS}=\SI{0.5}{\volt}$},
        што го задржува лебдечкото напојување на $\approx\SI{10.6}{\volt}$ --- далеку
        над \SI{6.7}{\volt} UVLO прагот на HB~\cite{ucc27211}.
  \item Полнеж од мирување/протекување во текот на најлошото време на спроводност на
        high-side-от. Со $I_{QBS}=\SI{0.12}{\milli\ampere}$~\cite{ucc27211}, MOSFET
        $I_{GSS}=2\times\SI{100}{\nano\ampere}$~\cite{ipt015ds} и висок duty cycle
        $D_{max}\approx 0.95$ (па $t_{HS}=0.95/\SI{40}{\kilo\hertz}=\SI{23.75}{\micro\second}$):
        \begin{equation}
        Q_{leak}\approx(\SI{0.12}{\milli\ampere})(\SI{23.75}{\micro\second})
                  \approx\SI{2.9}{\nano\coulomb}\ll Q_{g,\text{поз}}.
        \end{equation}
\end{itemize}
Полнежот на гејтот целосно доминира, па
\begin{equation}
Q_{TOTAL}\approx\SI{465}{\nano\coulomb}+\SI{3}{\nano\coulomb}\approx\SI{468}{\nano\coulomb},
\end{equation}
\begin{equation}
\boxed{\,C_{BOOT}\ge\dfrac{\SI{468}{\nano\coulomb}}{\SI{0.5}{\volt}}\approx\SI{0.94}{\micro\farad}\,}.
\end{equation}
Со примена на препорачаната резерва и со отчитување на опаѓањето на капацитивноста
при DC-преднапон кај керамичките кондензатори од класа\,II, избраната вредност е
\begin{quote}
\textbf{$C_{BOOT}=\SI{2.2}{\micro\farad}$, X7R, $\ge\SI{25}{\volt}$, керамички со
ниска ESR}, поставен директно меѓу пиновите HB--HS~\cite{ucc27211,infineon_maxboot}.
\end{quote}

\paragraph{Забелешка за опсегот \mbox{0.022--\SI{0.1}{\micro\farad}} од каталошкиот лист.}
Тој опсег е точен за \emph{еден} MOSFET со $\sim$\,30--\SI{90}{\nano\coulomb} полнеж
на гејтот~\cite{ucc27211}; овде паралелниот пар бара $\approx 5\times$ повеќе полнеж,
што е токму причината зошто каталошкиот лист наведува дека вредноста зависи од
полнежот на гејтот на high-side MOSFET-от~\cite{ucc27211}. Со употреба на
\SI{0.1}{\micro\farad} би се добило опаѓање од
$\SI{468}{\nano\coulomb}/\SI{0.1}{\micro\farad}\approx\SI{4.7}{\volt}$ по циклус ---
доволно за недоволно управување или дури за активирање на high-side UVLO. Затоа
поголемата вредност од \SI{2.2}{\micro\farad} не е опционална за оваа оптоварувачка
состојба.

\paragraph{Проверка на полнењето.} Кондензаторот го надополнува само
$\approx\SI{468}{\nano\coulomb}$ отстранет во секој циклус, а не целиот полнеж. При
најлош duty cycle, времето на спроводност на low-side-от е $\approx\SI{1.25}{\micro\second}$;
потребната просечна струја на полнење е
$\SI{468}{\nano\coulomb}/\SI{1.25}{\micro\second}\approx\SI{0.37}{\ampere}$, додека
интегрираната \SI{0.55}{\ohm} boot диода може почетно да испорача
$(12-0.9)/0.55\approx\SI{20}{\ampere}$~\cite{ucc27211}. Затоа полнењето е удобно при
\SI{40}{\kilo\hertz}. При многу мала брзина на моторот со duty cycle близу
$100\,\%$ на high-side-от, прозорецот за надополнување се намалува и може да биде
потребен периодичен low-side освежувачки импулс --- позната ограниченост на сите
bootstrap напојувања~\cite{infineon_monolithic,infineon_bootoverview}.

\subsection{Bootstrap диода}
UCC27211A-Q1 ја \emph{интегрира} bootstrap диодата (анода кон $V_{DD}$, катода кон
HB); таа е оценета на \SI{120}{\volt}, има типична динамичка отпорност од
\SI{0.55}{\ohm}, пад на напон од $\approx$\,0.9--\SI{1.05}{\volt} при
\SI{100}{\milli\ampere} и време на исклучување од \SI{20}{\nano\second}~\cite{ucc27211}.
За овој дизајн внатрешната диода е соодветна: Дел~\ref{sec:cboot} покажува дека и
струјата на полнење и опаѓањето се удобно во рамки на нејзината способност, а
нејзината \SI{120}{\volt} оценка ја надминува \SI{48}{\volt} сабирницата со голема
маргина~\cite{ucc27211}.

\emph{Опционална надворешна диода паралелно} со внатрешната може да се додаде ако
проектантот сака дополнително да го намали падот на напон при полнење и опаѓањето по
циклус. Ако се користи, мора да задоволи~\cite{infineon_monolithic,infineon_bootoverview}:
обратен напон $\ge V_{bus}+V_{DD}$ маргина $\rightarrow$ се избира \textbf{\SI{100}{\volt}}
компонента за совпаѓање со напонската класа на MOSFET-от; брзо обратно закрепнување
(reverse recovery, $t_{rr}\le\SI{50}{\nano\second}$) или Schottky; и струја на
спроводност $\ge$ врвната струја на полнење ($\ge\SI{1}{\ampere}$ е доволно). Бидејќи
функцијата веќе е обезбедена на чипот, оваа компонента е надградба, а не неопходност.

\subsection{Отпорници на гејтот и проверка на врвната струја / Miller}\label{sec:gateR}
\paragraph{Барање за пригушување (Balogh).} Оптималниот надворешен отпорник на гејтот
за критично пригушување на резонанцата во колото на гејтот, формирана од
индуктивноста на изворот $L_S$ и $C_{iss}$, е~\cite{balogh}
\begin{equation}
R_{GATE,OPT}=2\sqrt{\frac{L_S}{C_{ISS}}}-(R_{DRV}+R_{G,I}).
\end{equation}
За куќиштето TOLL со ниска индуктивност ($L_S$ од редот на \SI{1}{\nano\henry}) и
$C_{iss}=\SI{12}{\nano\farad}$, $2\sqrt{L_S/C_{ISS}}\approx\SI{0.58}{\ohm}$, додека
излезната импеданса на драјверот плус внатрешната отпорност на гејтот
($R_{DRV}+R_{G,I}$) веќе ја надминуваат. Резултатот е негативен, што значи дека
колото е \textbf{веќе прекумерно пригушено само од импедансата на драјверот и
внатрешната отпорност на гејтот од \SI{1.4}{\ohm}}~\cite{balogh,ipt015ds}. Затоа
надворешниот отпорник на гејтот се избира не за пригушување, туку за: (i) да ја држи
споделената излезна струја на драјверот во рамки на оценката, (ii) да го наштелува
$\mathrm{d}v/\mathrm{d}t$ заради EMI и (iii) да ги избалансира двата паралелни уреди.

\paragraph{Поединечни отпорници.} Најдобрата практика за паралелни MOSFET-и е
\emph{поединечен} отпорник на гејтот кај секој гејт (по два по прекинувачка позиција)
наместо еден заеднички, така што секоја тенденција двата уреди да осцилираат еден
против друг се пригушува и нивното прекинување се балансира. Се избира стандардна
вредност
\begin{quote}
\textbf{$R_{G,ext}=\SI{2.2}{\ohm}$ (по еден на секој MOSFET, E24)},
\end{quote}
блиску до \SI{1.8}{\ohm} при кои се карактеризирани времињата на прекинување на
IPT015N10N5~\cite{ipt015ds}.

\paragraph{Проверка на врвната струја / $\mathrm{d}v/\mathrm{d}t$.} Miller plateau е
на $V_{plateau}=\SI{4.4}{\volt}$~\cite{ipt015ds}. Со гејтот задржан на платото,
мрежата би \emph{барала} по уред
\begin{equation}
I_{G}=\frac{V_{DD}-V_{plateau}}{R_{G,ext}+R_{G,I}}=\frac{12-4.4}{2.2+1.4}\approx\SI{2.1}{\ampere},
\end{equation}
т.е.\ $\approx\SI{4.2}{\ampere}$ за двата уреди заедно. Ова ја надминува границата на
изворната струја на драјверот од \SI{3.7}{\ampere}~\cite{ucc27211}, па драјверот ја
ограничува струјата на \SI{3.7}{\ampere} и напонот на јазолот опаѓа --- драјверот, а
не отпорникот, ја одредува брзината. Времето за испорака на комбинираниот Miller
полнеж ($2\times\SI{51}{\nano\coulomb}=\SI{102}{\nano\coulomb}$ макс.) при
\SI{3.7}{\ampere} е
\begin{equation}
t_{Miller}=\frac{Q_{gd,\text{поз}}}{I_{O,source}}=\frac{\SI{102}{\nano\coulomb}}{\SI{3.7}{\ampere}}\approx\SI{28}{\nano\second},
\end{equation}
што дава, на најлошиот случај на сабирницата од \SI{60}{\volt} (полна \SI{48}{\volt} батерија),
\begin{equation}
\frac{\mathrm{d}V_{DS}}{\mathrm{d}t}\approx\frac{\SI{60}{\volt}}{\SI{28}{\nano\second}}\approx\SI{2.1}{\volt\per\nano\second},
\end{equation}
што е умерено (добро за спроведено EMI) и далеку под границата на slew rate од
\SI{50}{\volt\per\nano\second} на пинот HS~\cite{ucc27211}.

Ова го потврдува клучното прашање за компатибилноста: \textbf{дури и со удвоениот
полнеж на гејтот од паралелното поврзување, еден канал на UCC27211A-Q1 го прекинува
парот за под \SI{30}{\nano\second}} --- токму за што е наменета оценката
\SI{3.7}{\ampere}/\SI{4.5}{\ampere} на уредот, имено „управување на големи енергетски
MOSFET-и со минимизирани прекинувачки загуби при преминот низ Miller
plateau-то“~\cite{ucc27211}. Сепак, единствената вредност од \SI{2.2}{\ohm} изведена
овде е само \emph{симетрична основа}: Дел~\ref{sec:asym} ја префинува во асиметрична
мрежа за вклучување/исклучување, која е конфигурацијата што всушност е усвоена.

\subsection{Асиметрична (поделена) мрежа на отпорници на гејтот}\label{sec:asym}
Единствениот отпорник од Дел~\ref{sec:gateR} принудува една вредност да служи за
двата фронта на прекинување, што е компромис. Овој дел го заменува со
\textbf{асиметрична} мрежа што ги поставува брзините на вклучување и исклучување
независно --- побавно, помеко вклучување и брзо исклучување со ниска импеданса ---
која е конфигурацијата усвоена во финалниот дизајн.

\paragraph{Топологија.} По MOSFET, секогаш присутниот сериски отпорник
\textbf{$R_{on}$} ја носи струјата на вклучување (изворна), а гранка од
\textbf{$R_{off}$ во серија со OFF диода $D_{off}$} се поставува паралелно со него.
$D_{off}$ е ориентирана со анода кон гејтот и катода кон излезот на драјверот, така
што спроведува \emph{само кога гејтот се влече надолу} (исклучување /
$\mathrm{d}v/\mathrm{d}t$ стегач). Следствено,
\begin{equation}
R_{g(\text{on})}=R_{on},\qquad R_{g(\text{off})}=R_{on}\parallel(R_{off}+R_{D}).
\end{equation}
При вклучување $D_{off}$ е обратно поларизирана и гејтот гледа само $R_{on}$; при
исклучување двете патеки спроведуваат и гејтот ја гледа пониската паралелна
комбинација~\cite{balogh,pmeg4030er}. (Со една OFF диода, $R_{on}$ останува во
патеката на исклучување; пониската гранка $R_{off}$ едноставно го шунтира, поради што
$R_{g(off)}$ е паралелна комбинација, а не само $R_{off}$.)

\paragraph{Проектна насока --- бавно вклучување, брзо исклучување.} За тврдо
прекинуван полумост, претпочитаната асиметрија е \emph{побавно вклучување, побрзо
исклучување}~\cite{pmeg4030er,nexperia_gdfund,toshiba_dvdt}. Побрзото вклучување го
зголемува $\mathrm{d}i/\mathrm{d}t$ на струјата на одводот, што ја зголемува струјата
на обратно закрепнување (reverse recovery) на \emph{телесната freewheeling диода на
комплементарниот уред}; телесната диода на IPT015N10N5 складира
$Q_{rr}\approx\SI{316}{\nano\coulomb}$ ($t_{rr}\approx\SI{103}{\nano\second}$)~\cite{ipt015ds},
па намерно омекнатото вклучување го ограничува тој врв на закрепнување, придружната
загуба и EMI што таа ги генерира~\cite{pmeg4030er,nexperia_gdfund}. Спротивно на тоа,
исклучувањето има корист од тоа да биде брзо и со ниска импеданса: ги намалува
загубите при исклучување и, што е клучно, го држи гејтот на \emph{исклучениот} уред
цврсто кон масата, така што Miller струјата инјектирана низ $C_{GD}$ кога
комплементарниот прекинувач го менува јазолот не може да го подигне $V_{GS}$ до
прагот и да предизвика лажно (само-)вклучување~\cite{nexperia_gdfund,toshiba_dvdt}.
\SI{100}{\volt} уредот на сабирница до \SI{60}{\volt} има доволна напонска резерва
($\approx\SI{40}{\volt}$) да го апсорбира малку повисокиот $\mathrm{d}i/\mathrm{d}t$
надскок при исклучување што го произведува брзото исклучување; токму затоа што таа
резерва е помала отколку на фиксна \SI{48}{\volt} сабирница, можноста независно да се
контролира фронтот на исклучување е уште поценета. Асиметријата ја засилува и
сопствената нерамнотежа \SI{4.5}{\ampere}-понор / \SI{3.7}{\ampere}-извор на
драјверот~\cite{ucc27211}.

\paragraph{Избор на вредности.}
\begin{itemize}[leftmargin=1.4em]
  \item \textbf{$R_{on}=\SI{3.3}{\ohm}$ (по MOSFET).} Зголемен над компромисот од
        \SI{2.2}{\ohm} од Дел~\ref{sec:gateR} за да се омекне фронтот на вклучување.
        Резултирачката струја на вклучување по уред на платото е
        $I=(12-4.4)/(R_{on}+R_{G,I})=7.6/(3.3+1.4)=\SI{1.62}{\ampere}$, т.е.\
        $\approx\SI{3.2}{\ampere}$ за парот --- сега поставена од отпорникот, а не од
        границата на драјверот од \SI{3.7}{\ampere}, што дава контролиран
        $\mathrm{d}v/\mathrm{d}t$ од
        $\SI{60}{\volt}/(\SI{102}{\nano\coulomb}/\SI{3.2}{\ampere})\approx\SI{1.9}{\volt\per\nano\second}$
        во најлош случај~\cite{ipt015ds,ucc27211}.
  \item \textbf{$R_{off}=\SI{1.0}{\ohm}$ (по MOSFET).} Со динамичка отпорност на OFF
        диодата $R_{D}\approx$\,0.1--\SI{0.15}{\ohm} (проценета од спроводната
        карактеристика на PMEG4030ER, $V_F\approx$\,0.46--\SI{0.54}{\volt} во
        релевантниот опсег~\cite{pmeg4030er}):
        \begin{equation}
        R_{g(\text{off})}=R_{on}\parallel(R_{off}+R_D)=3.3\parallel 1.15\approx\SI{0.85}{\ohm}.
        \end{equation}
        Ова дава ефективен однос на отпорност вклучување:исклучување од околу
        \textbf{3.9\,:\,1}.
\end{itemize}

\paragraph{Зошто $R_{off}$ не се намалува понатаму.} \emph{Внатрешната} отпорност на
гејтот $R_{G,I}=\SI{1.4}{\ohm}$ е во серија со чипот и не може да се
заобиколи~\cite{ipt015ds}. Затоа вкупната отпорност на колото на гејтот при
исклучување е $R_{G,I}+R_{g(\text{off})}=1.4+0.85=\SI{2.25}{\ohm}$, и дури и со
$R_{off}=0$ не би можела да падне под $\approx\SI{1.4}{\ohm}$. Намалувањето на
$R_{off}$ под $\sim\SI{1}{\ohm}$ затоа дава сè помали придобивки, па \SI{1.0}{\ohm} е
добро поставен избор. При оваа вредност, врвната струја на исклучување по уред на
платото е $4.4/2.25\approx\SI{1.96}{\ampere}$ ($\approx\SI{3.9}{\ampere}$ за парот),
блиску до границата на понорната струја на драјверот од \SI{4.5}{\ampere}, што дава
транзиција на исклучување од $\approx\SI{26}{\nano\second}$ наспроти
$\approx\SI{32}{\nano\second}$ при вклучување --- посакуваната асиметрија~\cite{ucc27211}.

\paragraph{Избор на OFF диода.} $D_{off}$ мора: (i) да го блокира обратниот напон што
се појавува на $R_{on}$ при вклучување ($\approx I_{on,peak}\cdot R_{on}\approx
\SI{1.62}{\ampere}\times\SI{3.3}{\ohm}\approx\SI{5.3}{\volt}$, плус ringing); (ii) да
ја носи струјата на импулсот при исклучување ($\approx\SI{2}{\ampere}$ врв по гејт);
(iii) да закрепнува моментално, така што не додава полнеж на следното вклучување; и
(iv) да има ниска динамичка отпорност, така што $R_{off}$ доминира во патеката на
исклучување. Schottky диода ги задоволува сите четири. Избраната компонента е
\begin{quote}
\textbf{$D_{off}=$ Nexperia PMEG4030ER} --- \SI{40}{\volt}, \SI{3}{\ampere} Schottky
исправувачка диода со низок $V_{F}$ ($V_{F}\approx$\,0.46--\SI{0.54}{\volt},
$I_{FSM}=\SI{50}{\ampere}$, SOD-123W)~\cite{pmeg4030er}.
\end{quote}
Нејзината оценка од \SI{40}{\volt} дава $>7\times$ маргина над $\sim\SI{5}{\volt}$
обратен напрег; нејзината оценка од \SI{3}{\ampere} просечно / \SI{50}{\ampere}
ударно ја надминува струјата на импулсот при исклучување од $\sim\SI{2}{\ampere}$;
нејзиниот низок $V_F$ и мала спојна капацитивност ја држат патеката на исклучување
чиста; а како Schottky има занемарливо обратно закрепнување~\cite{pmeg4030er}. Широко
достапна алтернатива е 1N5819 (\SI{40}{\volt}, \SI{1}{\ampere} Schottky), со
прифаќање на нејзиното поголемо куќиште и капацитивност.

\paragraph{Снага и куќиште на отпорниците.} Просечната снага дисипирана во
отпорникот на гејтот е дел од вкупната загуба за управување на гејтот,
$P_G=Q_g\,V_{DD}\,f_{SW}$~\cite{balogh}. По уред, со $Q_g\approx\SI{232}{\nano\coulomb}$
при \SI{12}{\volt}, $P_G\approx\SI{232}{\nano\coulomb}\times\SI{12}{\volt}\times\SI{40}{\kilo\hertz}\approx\SI{0.11}{\watt}$,
поделена подеднакво меѓу вклучувањето и исклучувањето. Уделот што паѓа на надворешниот
отпорник е неговата фракција од вкупната отпорност во колото на гејтот~\cite{balogh};
со проценета излезна отпорност на драјверот $R_{HI}\approx\SI{3.2}{\ohm}$:
\begin{align}
P_{R_{on}} &\approx\frac{R_{on}}{R_{HI}+R_{on}+R_{G,I}}\cdot\tfrac{1}{2}Q_g V_{DD} f_{SW}
            =\frac{3.3}{7.9}\times\SI{0.0557}{\watt}\approx\SI{26}{\milli\watt},\\
P_{R_{off}} &\approx\SI{6}{\milli\watt}\ \text{(само при исклучување; останатото го носат $R_{HI}$ и $R_{G,I}$).}
\end{align}
Двете се далеку под $1/16$\,W. Сепак, отпорникот гледа кратки струјни врвови
($\sim$\,1.6--\SI{2}{\ampere}), па изборот на куќиште се води од енергијата по импулс
и од маргината, а не само од просечната снага: енергијата по транзиција во $R_{on}$
е $\approx 0.42\times\tfrac{1}{2}Q_g V_{DD}\approx\SI{0.6}{\micro\joule}$, занемарливо
во споредба со импулсната издржливост на кој било чип-отпорник. Се избира
\begin{quote}
\textbf{$R_{on}$, $R_{off}$: куќиште 0805, дебелослоен (thick-film), $1/8$\,W
(\SI{125}{\milli\watt})}, што дава $>4\times$ маргина на просечната снага и солидна
импулсна резерва; работниот напон ($\sim\SI{12}{\volt}$) е далеку под напонската
оценка на 0805 ($\sim\SI{150}{\volt}$)~\cite{balogh}.
\end{quote}
Куќиштето 0603 ($1/10$\,W) исто така е електрично доволно; 0805 е претпочитано заради
подобро ширење на топлината, импулсна робусност и леснотија на монтажа. (Pull-down
отпорникот $R_{GS}$ е димензиониран во Дел~\ref{sec:pulldown}.)

\begin{table}[t]
\centering
\caption{Споредба: единствен отпорник наспроти асиметрична мрежа.}
\label{tab:asym}
\small
\begin{tabularx}{\columnwidth}{@{}L L L@{}}
\toprule
\textbf{Аспект} & \textbf{Единствен $R_{G}=\SI{2.2}{\ohm}$ (Дел~\ref{sec:gateR})} &
\textbf{Асиметрична $R_{on}=\SI{3.3}{\ohm}$ / $R_{off}=\SI{1.0}{\ohm}$ + $D_{off}$} \\
\midrule
Фронт при вклучување & Среден --- принуден компромис & Помек ($R_{g(on)}=\SI{3.3}{\ohm}$);
$\mathrm{d}v/\mathrm{d}t\approx\SI{1.5}{\volt\per\nano\second}$~\cite{ipt015ds,ucc27211} \\
Фронт при исклучување & Ист како при вклучување (не може да се разликува) & Брз
($R_{g(off)}\approx\SI{0.85}{\ohm}$ надв.); близу полн понор од \SI{4.5}{\ampere}~\cite{ucc27211} \\
Обратно закрепнување на комплем.\ диода & Врзано за изборот при исклучување &
Независно укротено од бавното вклучување~\cite{pmeg4030er,nexperia_gdfund} \\
Импеданса на стегачот во исклучена состојба & $R_{G,I}+2.2=\SI{3.6}{\ohm}$ &
$R_{G,I}+0.85=\SI{2.25}{\ohm}\rightarrow$ повисок праг на $\mathrm{d}v/\mathrm{d}t$
за самовклучување~\cite{nexperia_gdfund,toshiba_dvdt} \\
Независна оптимизација на фронтовите & Невозможна & Да --- целта на мрежата~\cite{pmeg4030er} \\
Загуба при исклучување & Повисока ако $R_G$ се зголеми за меко вклучување & Пониска
(задржано брзо исклучување) \\
Дополнителни компоненти по MOSFET & --- & $+1$ отпорник, $+1$ Schottky \\
\bottomrule
\end{tabularx}
\end{table}

\paragraph{Зошто асиметричниот пристап е претпочитан.} Со единствен отпорник
проектантот се соочува со неизбежен конфликт~\cite{pmeg4030er,nexperia_gdfund}:
\emph{мала} вредност дава брзо исклучување (ниска загуба, силен стегач), но и брзо
вклучување што го влошува обратното закрепнување на комплементарната телесна диода,
$\mathrm{d}v/\mathrm{d}t$ и EMI; \emph{голема} вредност го омекнува вклучувањето, но
истовремено го забавува исклучувањето, зголемувајќи ја загубата при исклучување и
слабеејќи го стегачот на $\mathrm{d}v/\mathrm{d}t$ што спречува shoot-through.
Асиметричната мрежа го раствора овој конфликт со раздвојување на двата
фронта~\cite{pmeg4030er}. Тоа раздвојување е особено вредно во \emph{оваа}
апликација, бидејќи секоја од трите гранки е тврдо прекинуван полумост во кој секое
вклучување комутира freewheeling телесна диода што носи струја од редот на
\SI{100}{\ampere}; бидејќи два паралелни MOSFET-и со висок $Q_g$ по прекинувач ја
удвојуваат и Miller струјата и цената на секоја вкрстена спроводност; и бидејќи
напонската резерва \SI{100}{\volt}/\SI{48}{\volt} му дозволува на дизајнот да
искористи брзо исклучување без грижа за надскок. Цената --- дванаесет дополнителни
отпорници и дванаесет Schottky диоди за целиот инвертор --- е занемарлива наспроти
добивките во ефикасноста, EMI однесувањето и отпорноста на паразитско
вклучување~\cite{pmeg4030er,nexperia_gdfund,toshiba_dvdt}. Поради овие причини,
асиметричната мрежа ја заменува основата со единствен отпорник од Дел~\ref{sec:gateR}
и е конфигурацијата пренесена во финалниот список на материјали.

\subsection{Pull-down отпорници гејт--извор}\label{sec:pulldown}
Отпорник од гејт кон извор го држи секој MOSFET дефинирано исклучен пред излезите на
драјверот да станат активни при вклучување на напојувањето и обезбедува заштита од
Miller-индуцирано вклучување при $\mathrm{d}v/\mathrm{d}t$~\cite{balogh}. Вредноста на
pull-down-от следи од релацијата за имунитет на $\mathrm{d}v/\mathrm{d}t$~\cite{balogh}:
\begin{equation}
\left.\frac{\mathrm{d}v}{\mathrm{d}t}\right|_{LIMIT}=\frac{V_{TH}-0.007\,(T_J-25)}{R_{G,I}\,C_{GD}}.
\end{equation}
Во нормална работа, ниската излезна импеданса на драјверот и понорот од
\SI{4.5}{\ampere} доминираат во оваа заштита; главната задача на pull-down-от е
состојбата при вклучување на напојувањето. Стандардна вредност
\begin{quote}
\textbf{$R_{GS}=\SI{10}{\kilo\ohm}$ (по еден на секој MOSFET)}
\end{quote}
е доволно голема за да не го оптоварува драјверот и доволно мала за да го држи гејтот
надолу при вклучување. Проценката со капацитивен делител на индуцираното вклучување
ја потврдува маргината: кога спротивниот прекинувач наметнува најлош напон од
\SI{60}{\volt} на исклучениот уред, најлошиот пораст на гејтот низ делителот
$C_{rss}/C_{iss}=\SI{80}{\pico\farad}/\SI{12}{\nano\farad}\approx 0.0067$ би бил
$\approx\SI{0.40}{\volt}$ дури и со \emph{лебдечки} гејт --- сè уште далеку под прагот
од \SI{2.2}{\volt}~\cite{ipt015ds} --- а активниот понор од \SI{4.5}{\ampere} ја држи
вистинската вредност уште пониска~\cite{ucc27211}.

\paragraph{Снага и куќиште на $R_{GS}$.} Кога гејтот е управуван на \SI{12}{\volt},
$R_{GS}$ непрекинато дисипира $P=V_{GS}^2/R_{GS}=12^2/\SI{10}{\kilo\ohm}=\SI{14.4}{\milli\watt}$
(просечно помалку, бидејќи гејтот не е секогаш висок). Затоа е доволно куќиште
\textbf{0603, $1/10$\,W (\SI{100}{\milli\watt})} ($>6\times$ маргина); за
единственост на спецификацијата може да се користи и 0805, $1/8$\,W заедно со
отпорниците на гејтот.

\subsection{Кондензатори за декуплирање (decoupling)}
\paragraph{$V_{DD}$ (low-side) декуплирање.} Секогаш кога LO излезот испорачува струја
на гејтот, еднаков струен импулс се влече низ $V_{DD}$, па керамички кондензатор со
ниска ESR мора да стои директно меѓу $V_{DD}$--$V_{SS}$~\cite{ucc27211}.
Димензиониран за толерантно бранување $\Delta V$ со истиот аргумент за полнеж како кај
bootstrap-от,
\begin{equation}
C_{VDD}\ge\frac{Q_{g,\text{поз}}}{\Delta V}=\frac{\SI{465}{\nano\coulomb}}{\SI{0.5}{\volt}}\approx\SI{0.93}{\micro\farad}.
\end{equation}
TI препорачува \mbox{0.22--\SI{4.7}{\micro\farad}} и горниот крај за ниски температури
и висок полнеж на гејтот~\cite{ucc27211}. Избрано:
\begin{quote}
\textbf{$C_{VDD}=\SI{4.7}{\micro\farad}$ X7R $+$ \SI{0.1}{\micro\farad} керамички, по
драјвер}, со посебен кондензатор за секој од трите драјвери~\cite{ucc27211}.
\end{quote}

\paragraph{HB (high-side) декуплирање.} Оваа улога ја извршува самиот bootstrap
кондензатор (Дел~\ref{sec:cboot}), поставен директно меѓу пиновите HB--HS; посебната
препорака од TI за \mbox{0.022--\SI{0.1}{\micro\farad}} е \emph{минималното}
високофреквентно декуплирање и овде е надмината од пресметаната вредност од
\SI{2.2}{\micro\farad}~\cite{ucc27211}.

\paragraph{Bulk на \SI{12}{\volt} сабирницата.} Споделен bulk резервоар (на пр.\
\mbox{22--\SI{47}{\micro\farad}} електролитски или керамички) на помошната
\SI{12}{\volt} сабирница ги напојува трите драјвери и ги апсорбира збирните струјни
импулси за управување на гејтот; TI советува посебен $V_{DD}$ кондензатор за
декуплирање по драјвер во системи со повеќе драјвери~\cite{ucc27211}.

\paragraph{Декуплирање на DC-link.} Бидејќи изворот е батерија поврзана преку сноп
кабли, импедансата на изворот (внатрешната отпорност на батеријата плус
индуктивноста на снопот) е значителна; затоа локалното DC-link декуплирање е уште
поважно отколку кај крута лабораториска сабирница. Енергетскиот степен носи бранувачка
струја од редот на \SI{100}{\ampere} (до $\approx\SI{150}{\ampere}$ во најлош случај
при празна \SI{24}{\volt} батерија) и мора да се декуплира со DC-link банка со ниска
ESR/ESL (фолиски и/или керамички) близу полумостот, димензионирана за RMS бранувачката
струја и за да го држи надскокот на сабирницата во рамки на оценката од
\SI{100}{\volt} при комутација. Бидејќи највисокиот напон на батеријата е
\SI{60}{\volt}, кондензаторите на сабирницата треба да имаат напонска оценка
$\ge\SI{80}{\volt}$ (фолиски) или $\ge\SI{100}{\volt}$ (керамички, со отчитување на
опаѓањето при DC-преднапон). Прецизното димензионирање зависи од индуктивноста на
снопот и е задача на ниво на плоча.

\subsection{Список на компоненти за управување на гејтот (по гранка / по плоча)}
Табела~\ref{tab:gatebom} ги наведува пасивните помошни компоненти за управување на
гејтот. Покрај нив, активните уреди се \textbf{$12\times$ IPT015N10N5} (4 по гранка)
и \textbf{$3\times$ UCC27211A-Q1}.

\begin{table}[t]
\centering
\caption{Список на пасивни помошни компоненти за управување на гејтот.}
\label{tab:gatebom}
\small
\begin{tabularx}{\columnwidth}{@{}l L L c c@{}}
\toprule
\textbf{Ознака} & \textbf{Функција} & \textbf{Вредност} & \textbf{По гранка} & \textbf{3 гранки} \\
\midrule
$C_{BOOT}$ & Bootstrap кондензатор (HB--HS) & \SI{2.2}{\micro\farad} X7R $\ge\SI{25}{\volt}$ & 1 & 3 \\
$C_{VDD}$ & Декуплирање на $V_{DD}$ & \SI{4.7}{\micro\farad} X7R $+$ \SI{0.1}{\micro\farad} & 1 сет & 3 сета \\
$R_{on}$ & Отпорник за вклучување (по MOSFET) & \SI{3.3}{\ohm}, 0805, $1/8$\,W & 4 (2 HS $+$ 2 LS) & 12 \\
$R_{off}$ & Отпорник за исклучување (по MOSFET) & \SI{1.0}{\ohm}, 0805, $1/8$\,W & 4 & 12 \\
$D_{off}$ & OFF Schottky диода & Nexperia PMEG4030ER (\SI{40}{\volt}, \SI{3}{\ampere}) & 4 & 12 \\
$R_{GS}$ & Pull-down гејт--извор & \SI{10}{\kilo\ohm}, 0603, $1/10$\,W & 4 & 12 \\
$D_{BOOT}$ (опц.) & Надворешна паралелна boot диода & \SI{100}{\volt} брза/Schottky $\ge\SI{1}{\ampere}$ & 0--1 & 0--3 \\
$C_{BULK}$ & Bulk на \SI{12}{\volt} сабирница & \mbox{22--\SI{47}{\micro\farad}} & --- & 1 (споделен) \\
\bottomrule
\end{tabularx}
\end{table}

\section{Мерење на фазна струја}\label{sec:current}

\subsection{Преглед и споредба на топологии за мерење на струја}\label{sec:csurvey}
Постојат четири фамилии на мерење на струја кај инвертор; нивниот мерен јазол е
прикажан на Сл.~\ref{fig:topo} за еден half-bridge.

\begin{figure}[t]
\centering
\begin{circuitikz}[scale=0.92,transform shape,american]
% bus rails
\draw (0,4) node[anchor=east]{$\Vdcbus$} -- (6.5,4);
\draw (0,-1.2) node[anchor=east]{GND} -- (6.5,-1.2);
% high-side switch (box) with optional high-side shunt above
\draw (2,3.0) rectangle (3,3.8) node[midway]{$Q_H$};
\draw (2.5,4) -- (2.5,3.8);
\draw (2.5,3.0) -- (2.5,1.7);
% low-side switch
\draw (2,0.9) rectangle (3,1.7) node[midway]{$Q_L$};
\draw (2.5,1.7) -- (2.5,1.7);
% phase node
\draw (2.5,1.7) node[circ]{} -- (4.6,1.7) node[anchor=west]{фаза $\to$ мотор};
% inline shunt in phase
\draw (3.6,1.7) node[anchor=south]{\scriptsize inline $\Rshunt$};
\draw[fill=yellow!30] (3.25,1.55) rectangle (3.95,1.85);
% low-side shunt to ground
\draw (2.5,0.9) -- (2.5,0.2);
\draw[fill=yellow!30] (2.15,-0.1) rectangle (2.85,0.2) node[midway]{};
\draw (2.5,0.0) node[anchor=west]{\scriptsize low-side $\Rshunt$};
\draw (2.5,-0.1) -- (2.5,-1.2);
% high-side shunt position marker
\draw[fill=yellow!30] (2.15,4.0) rectangle (2.85,4.3);
\draw (2.5,4.45) node[anchor=south]{\scriptsize high-side $\Rshunt$ (сабирница)};
\draw (2.5,4.3) -- (2.5,5.0) -- (5.5,5.0) -- (5.5,4) ;
% DC-link single shunt
\draw[fill=orange!25] (5.15,-1.55) rectangle (5.85,-1.25);
\draw (5.5,-1.7) node[anchor=north]{\scriptsize единечен DC-link $\Rshunt$};
\draw (5.5,-1.2)--(5.5,-1.25); \draw(5.5,-1.55)--(5.5,-2.0)--(0,-2.0);
\end{circuitikz}
\caption{Можни позиции на shunt во една гранка од инверторот: high-side (на сериска
врска со гранката од сабирницата), low-side (на сериска врска со гранката кон
заземјувањето), in-line (на сериска врска со фазниот спроводник), и единечниот
DC-link shunt во заедничкиот повраток. Hall/магнетната алтернатива го заменува
shunt-от со сензор на поле околу фазниот спроводник.}
\label{fig:topo}
\end{figure}

\begin{table*}[t]
\centering
\caption{Споредба на топологии за мерење на струја кај трифазен инвертор.}
\label{tab:topo}
\small
\renewcommand{\arraystretch}{1.25}
\begin{tabular}{p{2.6cm} p{2.4cm} p{2.6cm} p{3.0cm} p{3.4cm}}
\toprule
\textbf{Атрибут} & \textbf{Low-side shunt} & \textbf{High-side shunt} & \textbf{In-line (in-phase) shunt} & \textbf{Магнетна / Hall} \\
\midrule
Common-mode напон на мерниот елемент & Близу до земја ($\approx\!0$) & На/близу $\Vdcbus$ & Се ниша од шина до шина на $\Fswsw$ со високо $\mathrm{d}v/\mathrm{d}t$ & Галвански изолиран; сензорот не гледа CM \\
Ја мери вистинската фазна струја при сите работни циклуси & Не --- само додека спроведува $Q_L$; слеп при $\approx100\%$ duty & Не --- само додека спроведува $Q_H$ & \textbf{Да --- континуирано, при кој било duty} & Да \\
Барање за засилувач & Low-CM CSA & High-CM CSA & High-CM CSA со PWM потиснување & --- (сензорот е засилувачот)\\
DC способност & Да & Да & Да & Да (Hall) / Не (CT, Rogowski)\\
Загуба од вметнување & $I^2R$ во shunt & $I^2R$ во shunt & $I^2R$ во shunt & Занемарлива\\
Точност / дрифт & Висока / низок & Висока / низок & Висока / низок & Умерена; офсет и термички дрифт, нелинеарност\\
Изолација & Не & Не & Не & Да\\
Релативна цена & Најниска & Ниска & Ниска--умерена (CSA/фаза) & Највисока\\
Најпогодна за & 6-step / евтини, ограничен duty опсег & Мониторинг на батерија/полнач & \textbf{Високи перформанси FOC низ целиот опсег на брзина/duty} & Многу високи струи или дизајни каде изолацијата е задолжителна\\
\bottomrule
\end{tabular}
\end{table*}

Решавачкиот атрибут за FOC е можноста да се семплира \emph{вистинската} фазна струја
во секој PWM период, независно од работниот циклус. Low-side shunt-от носи фазна
струја само додека спроведува low-side уредот, па при висок duty (висока брзина /
висока модулација) прозорецот на спроведување се намалува и мерењето станува
непрецизно или бара изобличување на PWM (вметнување на минимално време на
спроведување)~\cite{ti_lowside_threephase}. Единечниот DC-link shunt ги реконструира
трите фазни струи од моменталната склопна состојба, но трпи од „ненабљудливи“ области
во близина на границите на напонските вектори, каде што времињата на задржување на
два активни вектори паѓаат под времето на стабилизирање/blanking на shunt-от, што
повторно бара модификација на PWM. In-line (in-phase) shunt-от е поставен на сериска
врска со изводот на моторот, па ја спроведува фазната струја континуирано и може да
се семплира при кој било работен циклус; затоа токму оваа топологија се користи во
референтните дизајни за серво и за FOC со високи перформанси~\cite{tida00913,ina240ds}.
Нејзината единствена тешкотија е тоа што мерниот елемент се наоѓа на јазол што се ниша
помеѓу земја и $\Vdcbus$ на $\Fswsw$ со големо $\mathrm{d}v/\mathrm{d}t$, па
засилувачот мора да комбинира широк common-mode (CM) опсег со силно потиснување на CM
преодните појави~\cite{ina240ds,tida00913}.

Магнетната/Hall алтернатива ја отстранува загубата од вметнување и обезбедува
изолација, но по повисока цена и со поголем офсет и термички дрифт отколку ланецот
shunt-плус-zero-drift-засилувач; изолацијата не е потребна тука бидејќи дизајнот е
експлицитно common-ground. Со оглед на common-ground ограничувањето, потребата за FOC
набљудливост при целосен duty на сите три фази, и управливата дисипација во shunt-от
изведена подолу, избрана е топологијата \textbf{in-line shunt}, во согласност со
спецификацијата и со \mbox{48-V} in-line FOC референтниот дизајн на TI,
TIDA-00913~\cite{tida00913}.

\subsection{Димензионирање на shunt отпорникот}\label{sec:shunt}
Три конкурентни ограничувања ја одредуваат вредноста на shunt-от $\Rshunt$: (1)
дисипацијата на моќност $P_{\Rshunt}=\Iphrms^2\Rshunt$ мора да остане во рамки на
економичен површински-монтиран елемент; (2) сигналниот напон на полна скала
$\hat V_{\mathrm{sh}}=\hat I_{\mathrm{ph}}\Rshunt$ мора, по засилувањето, да се
вклопи во влезниот прозорец на ADC без отсекување (clipping) при врвот во најлош
случај (со резерва за преоптоварување); и (3) $\Rshunt$ треба да биде доволно голем
за да офсетот и шумот на засилувачот, сведени на струја, останат мали.

За $\Rshunt=\SI{0.1}{\milli\ohm}$ при работната точка во најлош случај
($\Iphrms^{\max}=\SI{136}{\ampere}$ rms):
\begin{align}
P_{\Rshunt} &=\Iphrms^{2}\Rshunt=(136.1)^2\times 0.1\times10^{-3}=\SI{1.85}{\watt},\\
\hat V_{\mathrm{sh}} &=\hat I_{\mathrm{ph}}\,\Rshunt=192.5\times 0.1\times10^{-3}=\SI{19.3}{\milli\volt}.
\end{align}
Отпор од \SI{0.1}{\milli\ohm} ја држи дисипацијата на \SI{1.85}{\watt}, што удобно се
справува од страна на shunt од метален елемент (метална лента) со \SI{5}{\watt}
(резерва од $2.7\times$), додека \SI{19.3}{\milli\volt} врв е доволно голем за
прецизно засилување. Елемент од метална лента е избран поради неговата многу ниска
паразитна индуктивност (ESL) и низок температурен коефициент; четири-приклучна
(Kelvin) врска е задолжителна за сигналниот напон да се зема само на резистивниот
елемент, а не на лемните споеви со голема струја~\cite{ina240ds}. ESL е важна бидејќи
измерениот напон е $v_{\mathrm{sh}}(t)=\Rshunt\,i(t)+\mathrm{ESL}\,\dfrac{\mathrm{d}i}{\mathrm{d}t}$;
на \SI{40}{\kilo\hertz} со брзи склопни рабови индуктивниот член мора да биде
занемарлив наспроти резистивниот член од \SI{0.1}{\milli\ohm}, што го постигнува
елемент од метална лента со $<\SI{2}{\nano\henry}$.

\subsection{Избор на струен засилувач и засилување}\label{sec:csa}
Засилувачот мора (i) да поднесува CM напон што се ниша од малку под земја
(спроведување на body-диода / dead-time, $\approx\SI{-1}{\volt}$) до
$\Vdcbus=\SI{60}{\volt}$ плус пречекорувањето од ѕвонење; (ii) да го потиснува
големото CM $\mathrm{d}v/\mathrm{d}t$ при секој склопен раб така што излезот да не
ѕвони; (iii) да биде \emph{двонасочен} бидејќи фазната струја е AC; и (iv) да има
пропусен опсег далеку над \SI{40}{\kilo\hertz}. Texas Instruments INA240 ги задоволува
сите четири: тоа е струен засилувач со напонски излез и zero-drift архитектура, со CM
опсег од $-4$ до $+\SI{80}{\volt}$ \emph{независно од напојувањето}, со „enhanced PWM
rejection“ експлицитно проектирано да ги потисне CM преодните појави на PWM за погон
на мотор, и се нуди во четири фиксни засилувања од $20,50,100,200$~V/V~\cite{ina240ds}.
Долната граница од $-\SI{4}{\volt}$ ја прифаќа екскурзијата под земја, а горната
граница од $+\SI{80}{\volt}$ остава \SI{20}{\volt} резерва над сабирницата од
\SI{60}{\volt} за пречекорувањето од ѕвонење.

При напојување од \SI{3.3}{\volt} аналоген извор, двонасочниот излез при нулта струја
е центриран од двата референтни приклучоци на INA240, чиј внатрешен разделник ја
поставува средната точка на излезот на
$V_{\mathrm{ref}}=(V_{\mathrm{REF1}}+V_{\mathrm{REF2}})/2$~\cite{ina240ds}; врзувањето
на $\mathrm{REF1}$ за земја и $\mathrm{REF2}$ за шината од \SI{3.3}{\volt} дава
$V_{\mathrm{ref}}=\SI{1.65}{\volt}$. Тогаш излезниот напон е
\begin{equation}
V_{\mathrm{out}}=V_{\mathrm{ref}}+G\,\Rshunt\,i_{\mathrm{ph}} .
\label{eq:csaout}
\end{equation}
Со rail-to-rail излез, употребливата едностранска амплитуда е $\approx\SI{1.45}{\volt}$
(со резерва од \SI{0.2}{\volt} до секоја шина). Со избор на засилување $G=50$
(INA240A2), линеарната врвна струја на полна скала е
\begin{equation}
\hat I_{\mathrm{FS}}=\frac{1.45}{G\,\Rshunt}
   =\frac{1.45}{50\times 0.1\times10^{-3}}=\SI{290}{\ampere}\ \text{врв}\;(\SI{205}{\ampere}\,\text{rms}).
\end{equation}
Врвот во најлош случај од \SI{193}{\ampere} произведува излезна екскурзија од
$G\Rshunt\hat I_{\mathrm{ph}}=50\times0.1\text{m}\times192.5=\SI{0.96}{\volt}$, т.е.\
$V_{\mathrm{out}}\in[0.69,2.61]\,\si{\volt}$, оставајќи $\approx\SI{0.49}{\volt}$ пред
отсекувањето. Оваа резерва е намерно зачувана за преодно пречекорување и овозможува
прагот за пречекорна струја (over-current, OC) да се постави на $\approx\SI{250}{\ampere}$
врв ($1.3\times\hat I_{\mathrm{ph}}^{\max}$, излезна екскурзија \SI{1.25}{\volt}), што
може да се детектира од аналоген компаратор на STM32G4 или софтверски пред засилувачот
да засити. Квантизацијата и офсетот, сведени на струја, се
\begin{align}
\Delta I_{\mathrm{LSB}} &=\frac{V_{\mathrm{ref+}}/2^{N}}{G\,\Rshunt}
   =\frac{3.3/4096}{50\times0.1\text{m}}=\SI{0.16}{\ampere/LSB},\\
I_{\mathrm{os}} &=\frac{V_{\mathrm{os(max)}}}{\Rshunt}=\frac{\SI{25}{\micro\volt}}{0.1\text{m}\Omega}
   =\SI{0.25}{\ampere},
\end{align}
двете занемарливи наспроти сигнал од \SI{136}{\ampere} (zero-drift архитектурата на
INA240 дава $V_{\mathrm{os}}\le\SI{25}{\micro\volt}$~\cite{ina240ds}). Поради тоа,
парот засилување-50 / \SI{0.1}{\milli\ohm} се претпочита пред алтернативниот пар
засилување-20 / \SI{0.25}{\milli\ohm}, кој би дисипирал \SI{4.6}{\watt} по shunt за
иста полна скала.

\subsection{ADC интерфејс и излезно филтрирање}\label{sec:csafilter}
За FOC фазната струја се семплира \emph{синхроно} со PWM носителот (вообичаено при
врвот/долот на бројачот), така што моментот на семплирање се совпаѓа со средината на
склопната состојба, каде што компонентата на бранувањето (ripple) на индукторот е на
нејзината средна вредност; ова природно го потиснува бранувањето на \SI{40}{\kilo\hertz}
и неговите хармоници без тежок филтер~\cite{an5346}. Тежок anti-alias филтер наместо
тоа би додал фазно доцнење на токму оној сигнал на кој се затвора струјната јамка,
влошувајќи ја јамката. Затоа излезот на CSA користи само лесен RC филтер од прв ред
($R_f=\SI{100}{\ohm}$, $C_f=\SI{2.2}{\nano\farad}$), кој (i) го ограничува опсегот на
високофреквентниот склопен шум и ѕвонење (гранична фреквенција
$f_c=1/2\pi R_fC_f=\SI{723}{\kilo\hertz}$, далеку над содржината на струјната јамка во
kHz-опсегот, но придушувајќи го ѕвонењето од десетици MHz), и (ii) дејствува како
резервоар на полнеж за sample-and-hold на ADC. Ниската излезна импеданса на INA240 ги
управува оваа мрежа и кондензаторот за семплирање на ADC со обилна резерва за
стабилизирање. Хардверското oversampling/усреднување на STM32G4 може да се овозможи за
дополнително намалување на шумот по цена на пропусност~\cite{an5346}.
Сл.~\ref{fig:csa} го прикажува комплетниот ланец по фаза.

\begin{figure}[t]
\centering
\begin{circuitikz}[scale=1.0,transform shape,american]
% shunt row (top)
\draw (0,4.7) node[anchor=east]{фазен јазол} -- (1.3,4.7);
\draw[fill=yellow!30] (1.3,4.5) rectangle (2.5,4.9);
\draw (1.9,5.2) node{\scriptsize $R_{\mathrm{sh}}=\SI{0.1}{\milli\ohm}$ (Kelvin)};
\draw (2.5,4.7) -- (3.9,4.7) node[anchor=west]{кон мотор};
% Kelvin taps (non-crossing): right terminal -> IN+ (upper), left terminal -> IN- (lower)
\draw (2.3,4.5) -- (2.3,3.0) -- (2.7,3.0);
\draw (1.5,4.5) -- (1.5,2.5) -- (2.7,2.5);
\draw (2.45,3.0) node[anchor=south]{\scriptsize IN+};
\draw (2.0,2.5) node[anchor=south]{\scriptsize IN--};
% INA240 block
\draw (2.7,1.9) rectangle (5.1,3.5);
\draw (3.9,2.45) node{\small INA240A2};
\draw (3.9,3.05) node{\scriptsize $G=50$};
% supply stub (top, label kept clear of the shunt row)
\draw (3.4,3.5) -- (3.4,4.0) node[anchor=east]{\scriptsize $3V3_{\mathrm A}$+\SI{100}{\nano\farad}};
% reference stub (bottom)
\draw (3.9,1.9) -- (3.9,1.2) node[anchor=north,align=center]{\scriptsize REF1=GND, REF2=3V3\\\scriptsize $\Rightarrow V_{\mathrm{ref}}=\SI{1.65}{\volt}$};
% output + filter (right)
\draw (5.1,2.7) -- (5.6,2.7) to[R,l={$R_f$, \SI{100}{\ohm}}] (7.1,2.7);
\draw (7.1,2.7) node[circ]{} to[C,l={$C_f$, \SI{2.2}{\nano\farad}}] (7.1,1.5) node[ground]{};
\draw (7.1,2.7) -- (7.8,2.7) node[anchor=west]{\scriptsize ADC};
\end{circuitikz}
\caption{In-line ланец за мерење на струја по фаза: Kelvin-приклучен shunt од метална
лента со \SI{0.1}{\milli\ohm}, INA240A2 ($G=50$) со референца на средна скала, и лесен
RC од \SI{723}{\kilo\hertz} во ADC на STM32G4. Се користат три идентични ланци (фази
U, V, W).}
\label{fig:csa}
\end{figure}

\section{Мерење на напон}\label{sec:voltage}

\subsection{Барања}
Контролерот мора да поддржува сензорско FOC (Hall/енкодер) и безсензорско FOC, при
што второто ја реконструира положбата на роторот од back-EMF. Затоа е потребно
хардверско мерење на напон за (i) напонот на DC-сабирницата --- за заштита од
пренапон/поднапон, скалирање на индексот на модулација и пресметка на моќност --- и
(ii) секој од трите напони на фазните јазли --- за набљудување на back-EMF и за
повратна врска по фазен напон. Сите четири канали го делат заземјувањето на инверторот
(неизолиран дизајн), па секој е заземјен отпорнички разделник што напојува ADC канал на
STM32G4. Каналот на сабирницата е спор (заштита/скалирање), додека фазните канали мора
да го зачуваат тајмингот на back-EMF за набљудувачот на положбата.

\subsection{Дизајн на отпорнички разделник}\label{sec:divider}
Опсегот на конверзија на ADC е $0$--$V_{\mathrm{REF+}}$. Користењето
$V_{\mathrm{REF+}}=V_{\mathrm{DDA}}=\SI{3.3}{\volt}$ го максимизира опсегот; сооднос
на разделникот $k=(R_{\mathrm{top}}+R_{\mathrm{bot}})/R_{\mathrm{bot}}$ го пресликува
максималниот измерен напон во близина на полна скала, со резерва. Со избор на
$R_{\mathrm{top}}=\SI{110}{\kilo\ohm}$ и $R_{\mathrm{bot}}=\SI{4.7}{\kilo\ohm}$,
\begin{equation}
k=\frac{R_{\mathrm{top}}+R_{\mathrm{bot}}}{R_{\mathrm{bot}}}
 =\frac{110+4.7}{4.7}=24.4,\qquad V_{\mathrm{ADC}}=\frac{V_{\mathrm{in}}}{k}.
\end{equation}
Ова пресликува $\Vdcbus=\SI{60}{\volt}\to\SI{2.46}{\volt}$ (\SI{73.2}{\volt} при
полна скала, т.е.\ \SI{22}{\percent} резерва над опсегот) и
$\SI{30}{\volt}\to\SI{1.23}{\volt}$. Истиот сооднос се користи за фазните разделници
бидејќи фазен јазол може да го достигне напонот на сабирницата. Постојаната струја и
дисипацијата при \SI{60}{\volt} се
\begin{equation}
I_{\mathrm{div}}=\frac{60}{R_{\mathrm{top}}+R_{\mathrm{bot}}}=\SI{523}{\micro\ampere},\quad
P_{R_{\mathrm{top}}}=I_{\mathrm{div}}^2R_{\mathrm{top}}=\SI{30}{\milli\watt},
\end{equation}
па сите четири разделници заедно повлекуваат $\approx\SI{126}{\milli\watt}$ ---
занемарливо. Горниот отпорник го гледа целиот напон на сабирницата; единечен
thick-film чип од \SI{110}{\kilo\ohm} со номинал $\ge\SI{150}{\volt}$ (на пр.\
0805/1206) е адекватен при \SI{60}{\volt}, но се препорачува делење на
$R_{\mathrm{top}}$ на два сериски делови од \SI{55}{\kilo\ohm} за да се преполови
напонскиот стрес по компонента и да се подобри температурното совпаѓање на соодносот.
Th\'evenin изворната импеданција претставена на ADC е
\begin{equation}
R_{\mathrm{src}}=R_{\mathrm{top}}\parallel R_{\mathrm{bot}}
   =\frac{110\text{k}\times4.7\text{k}}{110\text{k}+4.7\text{k}}=\SI{4.51}{\kilo\ohm},
\label{eq:rsrc}
\end{equation}
што е клучната големина за прашањето за бафер во Дел~\ref{sec:buffer}.

\subsection{Дали ADC бара op-amp бафер? Квантитативен одговор}\label{sec:buffer}
ADC канал на STM32G4 семплира така што го поврзува изворот, преку внатрешна отпорност
на мултиплексер/прекинувач $R_{\mathrm{ADC}}$, на кондензатор за семплирање
$C_{\mathrm{ADC}}$ во текот на програмираното време на семплирање $t_s$
(Сл.~\ref{fig:adcmodel}). Кондензаторот мора да се стабилизира во рамки на
$\tfrac12\,$LSB од изворниот напон. За $N$-битна конверзија ова бара~\cite{an2834,an5346}
\begin{equation}
t_s\;\ge\;\bigl(R_{\mathrm{AIN}}+R_{\mathrm{ADC}}\bigr)\,C_{\mathrm{ADC}}\,\ln\!\bigl(2^{\,N+1}\bigr),
\label{eq:settle}
\end{equation}
каде што $R_{\mathrm{AIN}}=R_{\mathrm{src}}$ е надворешната изворна импеданција. ST не
ги објавува $R_{\mathrm{ADC}}$/$C_{\mathrm{ADC}}$ директно за G4; наместо тоа
техничкиот лист дава табела за „максимален $R_{\mathrm{AIN}}$ наспроти време на
семплирање“, а упатството на ST е да се избере времето на семплирање од таа табела за
дадената изворна импеданција~\cite{an5346}. Со користење на моделот на стабилизирање
од AN2834 со репрезентативни вредности $C_{\mathrm{ADC}}\approx\SI{5}{\pico\farad}$ и
$R_{\mathrm{ADC}}\approx\SI{2}{\kilo\ohm}$, и $\ln(2^{13})=9.01$ за 12-битниот режим,
\begin{equation}
t_s\ge(4510+2000)(5\times10^{-12})(9.01)=\SI{293}{\nano\second}.
\end{equation}
При ADC такт од \SI{60}{\mega\hertz} ова е $17.6$ ADC циклуси, додека STM32G4 нуди
поставки за време на семплирање до $640.5$ циклуси; самата поставка од $47.5$ циклуси
е \SI{792}{\nano\second} --- повеќе од $2.7\times$ од побараното. Бидејќи напоните на
сабирницата и фазите бараат само освежување со kHz-стапка, потрошокот на подолго време
на семплирање не чини ништо. Затоа:
\begin{center}
\emph{Не е потребен надворешен op-amp бафер за напонските разделници.}
\end{center}
Изворната импеданција на разделникот се полни до 12-битна точност во рамки на време на
семплирање што е тривијално достапно; мал филтерски кондензатор (Дел~\ref{sec:vfilter})
дополнително ја намалува ефективната изворна импеданција во моментот на семплирање.

\begin{figure}[t]
\centering
\begin{circuitikz}[scale=1.0,transform shape,american]
\draw (0,1.2) node[anchor=east]{$V_{\mathrm{div}}$} to[R,l={$R_{\mathrm{AIN}}{=}R_{\mathrm{src}}$}] (3,1.2);
\draw (3,1.2) to[R,l=$R_{\mathrm{ADC}}$] (5,1.2);
\draw (5,1.2) to[nos,l=$t_s$] (6.2,1.2);
\draw (6.2,1.2) to[C,l=$C_{\mathrm{ADC}}$] (6.2,0) node[ground]{};
\draw (6.2,1.2) -- (7.0,1.2) node[anchor=west]{SAR};
\end{circuitikz}
\caption{Модел од прв ред на влезниот степен за семплирање на ADC на STM32G4 користен
во~\eqref{eq:settle}: изворната импеданција на разделникот $R_{\mathrm{src}}$ на
сериска врска со внатрешната отпорност на прекинувачот го полни кондензаторот за
семплирање $C_{\mathrm{ADC}}$ во текот на прозорецот на семплирање $t_s$.}
\label{fig:adcmodel}
\end{figure}

\begin{table}[t]
\centering
\caption{Op-amp бафер пред ADC: придобивки наспроти трошоци.}
\label{tab:buffer}
\small
\renewcommand{\arraystretch}{1.2}
\begin{tabular}{p{4.0cm} p{4.0cm}}
\toprule
\textbf{Аргументи за бафер} & \textbf{Аргументи против бафер}\\
\midrule
Ниска излезна импеданса $\Rightarrow$ брзо стабилизирање, кратко време на семплирање, висока стапка на мултиплексирано скенирање &
Додава сопствен офсет, дрифт, шум и грешка од ограничен пропусен опсег на \emph{апсолутно} мерење\\
Го изолира разделникот од „kickback“ поради инјектирање на полнеж во ADC / преслушување меѓу скенираните канали &
Дополнителен трошок на BOM, површина на плочата и декуплирање на напојувањето\\
Овозможува активен anti-alias филтер &
Станува дополнителен товар за калибрација\\
Го спречува резистивното оптоварување на разделникот &
Непотребен тука: бараното $t_s$ е мало,~\eqref{eq:settle}\\
\bottomrule
\end{tabular}
\end{table}

Ако некогаш би се посакувал бафер --- на пример за делење на еден разделник меѓу
повеќе брзи мултиплексирани канали, или за додавање активен филтер --- тој треба да се
реализира со \emph{сопствените ресурси на MCU} наместо со надворешен IC: STM32G473
содржи шест операционни засилувачи што можат да се конфигурираат како follower-и со
единечно засилување (или PGA) со нивниот излез насочен \emph{внатрешно} кон ADC канал,
што истовремено го отстранува патот на собирање шум преку надворешен приклучок и не
чини ништо на BOM~\cite{an5306,stm32g473ds}. ST токму оваа follower-кон-ADC употреба ја
документира и препорачува време на семплирање од \SI{200}{\nano\second} при читање на
излезот од внатрешен op-amp~\cite{an5346,an5306}. Компромисот околу баферот е сумиран
во Табела~\ref{tab:buffer}; за тековните напонски канали со ниска стапка заклучокот е
да се изостави баферот и да се користи подолго време на семплирање плюс филтерски
кондензатор.

\subsection{Филтрирање на сабирницата и фазите, и заштита}\label{sec:vfilter}
Кондензатор $C_{\mathrm{f}}$ паралелно со $R_{\mathrm{bot}}$ формира нископропусен
филтер од прв ред со гранична фреквенција
\begin{equation}
f_c=\frac{1}{2\pi R_{\mathrm{src}}C_{\mathrm{f}}} .
\end{equation}
\textbf{Канал на сабирницата:} $C_{\mathrm{f}}=\SI{22}{\nano\farad}$ дава
$f_c=\SI{1.6}{\kilo\hertz}$, што го отстранува склопното бранување (ripple) и шумот од
бавното мерење на сабирницата. \textbf{Фазни канали (back-EMF):} филтерот мора да го
отфрли PWM од \SI{40}{\kilo\hertz} додека ја пропушта основната хармоника на back-EMF.
Електричната фреквенција е $f_e=\mathrm{RPM}\cdot p/60$ (со $p$ полови парови); изборот
на $f_c\approx$\,2--\SI{3}{\kilo\hertz} го придушува носителот додека ја пропушта
основната хармоника, по цена на фазно доцнење $\phi=\arctan(f_e/f_c)$.
Табела~\ref{tab:phase} го наведува ова доцнење; тоа е детерминистичко и се компензира
во firmware (набљудувачот ја поместува семплираната фаза за $\phi$). За безсензорска
комутација базирана на премин низ нула, истото фиксно доцнење се коригира како
временски офсет.

\begin{table}[t]
\centering
\caption{Доцнење на нископропусниот фазен филтер $\phi=\arctan(f_e/f_c)$ (компензирано во firmware).}
\label{tab:phase}
\small
\begin{tabular}{c c c c}
\toprule
$f_c$ & $f_e=\SI{200}{\hertz}$ & $f_e=\SI{500}{\hertz}$ & $f_e=\SI{1000}{\hertz}$\\
\midrule
\SI{2}{\kilo\hertz} & \ang{5.7} & \ang{14.0} & \ang{26.6}\\
\SI{3}{\kilo\hertz} & \ang{3.8} & \ang{9.5}  & \ang{18.4}\\
\bottomrule
\end{tabular}
\end{table}

\textbf{Заштита:} за време на dead-time/спроведување на body-диода фазен јазол се ниша
до $\approx\SI{-1}{\volt}$, што разделникот го скалира на
$-1/24.4=\SI{-41}{\milli\volt}$ на приклучокот на ADC --- мало, но Schottky ограничувач
(clamp, на пр.\ BAT54S) од јазолот на ADC кон земја и кон $V_{\mathrm{DDA}}$ е додаден
како осигурување против преодни појави над/под опсегот, при што долниот отпорник на
разделникот ја ограничува струјата на ограничувачот. Истиот ограничувач штити од врвови
на пренапон на каналот на сабирницата. Сл.~\ref{fig:divider} ги прикажува разделникот,
филтерот и ограничувачот.

\begin{figure}[t]
\centering
\begin{circuitikz}[scale=1.05,transform shape,american]
% input and top resistor
\draw (0,4.2) node[anchor=south]{$V_{\mathrm{in}}$ (сабирница / фаза)} -- (0,3.4);
\draw (0,3.4) to[R,l={$R_{\mathrm{top}}$, \SI{110}{\kilo\ohm}}] (0,1.7);
\draw (0,1.7) node[circ]{};
% bottom resistor to ground
\draw (0,1.7) to[R,l={$R_{\mathrm{bot}}$, \SI{4.7}{\kilo\ohm}}] (0,0) node[ground]{};
% horizontal sense line to ADC
\draw (0,1.7) -- (5.2,1.7);
% filter cap tap
\draw (2.3,1.7) node[circ]{} to[C,l=$C_f$] (2.3,0) node[ground]{};
% clamp diodes (ESD pair) at x=3.4
\draw (3.4,1.7) node[circ]{};
\draw (3.4,1.7) to[Do] (3.4,3.4) node[anchor=south]{$V_{\mathrm{DDA}}$};
\draw (3.4,0) node[ground]{} to[Do] (3.4,1.7);
\draw (4.3,3.0) node[anchor=west]{\scriptsize BAT54S};
\draw (5.2,1.7) node[anchor=west]{кон ADC};
\end{circuitikz}
\caption{Напонски разделник за мерење (сооднос $24.4$), филтерски кондензатор од прв
ред, и Schottky ограничувач кон $V_{\mathrm{DDA}}$/GND. Идентични мрежи ги опслужуваат
DC-сабирницата ($f_c\approx\SI{1.6}{\kilo\hertz}$) и трите фазни јазли
($f_c\approx$\,2--\SI{3}{\kilo\hertz} за back-EMF).}
\label{fig:divider}
\end{figure}

\section{RC snubber на склопниот јазол}\label{sec:snubber}

\subsection{Потекло на ѕвонењето (ringing)}
При секое исклучување, паразитната индуктивност $L_p$ на комутационата јамка (изводи
на пакувањето, PCB траги, ESL на кондензаторот на сабирницата) резонира со излезната
капацитивност на MOSFET $C_o$ ($C_{\mathrm{oss}}$ на уредот плюс монтажната
капацитивност), произведувајќи придушена осцилација на склопниот јазол на
\begin{equation}
f_{\mathrm{ring}}=\frac{1}{2\pi\sqrt{L_p C_o}} ,
\label{eq:fring}
\end{equation}
вообичаено во десетици MHz. Ѕвонењето го оптоварува напонот на уредот, зрачи EMI и се
спрегнува во сигналните линии. RC придушувачки snubber поставен паралелно со уредот
додава отпорник $R_s$ (на сериска врска со кондензатор $C_s$) кој ја дисипира
резонантната енергија~\cite{severns,digikey_snubber}.

\subsection{Методологија на дизајн}
Ригорозната постапка, независна од layout-от, го мери ѕвонењето на склопената
плоча~\cite{severns,digikey_snubber}:
\begin{enumerate}[leftmargin=1.4em]
\item Измери го периодот на ѕвонење без snubber $T_1=1/f_{\mathrm{ring}}$ на склопниот јазол.
\item Додај познат кондензатор $C_{\mathrm{test}}$ паралелно со уредот и измери го новиот период $T_2$.
\item Бидејќи $T_1=2\pi\sqrt{L_pC_o}$ и $T_2=2\pi\sqrt{L_p(C_o+C_{\mathrm{test}})}$, квадрирањето и
одземањето го елиминира $L_p$:
\begin{equation}
L_p=\frac{T_2^{2}-T_1^{2}}{4\pi^{2}\,C_{\mathrm{test}}},\qquad
C_o=\frac{T_1^{2}}{4\pi^{2}L_p}.
\label{eq:extract}
\end{equation}
\end{enumerate}
Придушувачкиот отпорник се поставува на карактеристичната импеданција на резонанцата,
\begin{equation}
R_s=\sqrt{\frac{L_p}{C_o}},
\label{eq:rs}
\end{equation}
што дава приближно критично придушување~\cite{digikey_snubber}. Snubber кондензаторот
се избира неколку пати поголем од $C_o$ така што $R_s$ да доминира во јамката, а сепак
доволно мал за да се ограничи дисипацијата; моќноста на отпорникот е енергијата
складирана во $C_s$, полнета и празнета еднаш по склопен циклус,
\begin{equation}
P_{R_s}=C_s\,\Vdcbus^{2}\,\Fswsw .
\label{eq:psnub}
\end{equation}

\subsection{Проценети вредности (да се валидираат на тест-маса)}
Сè уште не постои склопена плоча, па $L_p$ не може да се измери; вредностите подолу се
прва проценка само за избор на footprint и \emph{мора} да се потврдат со методот со
две мерења од~\eqref{eq:extract} на прототипот. Од техничкиот лист на IPT015N10N5,
$C_{\mathrm{oss}}=\SI{1800}{\pico\farad}$ (наведено при тест-условот од техничкиот лист;
$C_{\mathrm{oss}}$ силно зависи од напонот)~\cite{ipt015ds,ipt015part}; додавањето
$\approx\SI{200}{\pico\farad}$ монтажна капацитивност дава $C_o\approx\SI{2}{\nano\farad}$.
Земајќи репрезентативно $L_p\approx\SI{10}{\nano\henry}$ за збиен layout,
\begin{align}
f_{\mathrm{ring}}&=\frac{1}{2\pi\sqrt{(10\text{n})(2\text{n})}}=\SI{35.6}{\mega\hertz},\\
R_s&=\sqrt{\frac{10\text{nH}}{2\text{nF}}}=\SI{2.24}{\ohm}\;\to\;\SI{2.2}{\ohm}\ (\text{E12}).
\end{align}
Со избор на $C_s=\SI{2.2}{\nano\farad}$ ($\approx C_o$), дисипацијата во најлош случај
(сабирница од \SI{60}{\volt}) од~\eqref{eq:psnub} е
\begin{equation}
P_{R_s}=2.2\text{n}\times 60^{2}\times 40\text{k}=\SI{0.32}{\watt}\ \text{по snubber},
\end{equation}
т.е.\ $\le\SI{1.9}{\watt}$ за сите шест прекинувачи ако секој footprint е монтиран
(само \SI{0.08}{\watt} по секој на сабирницата од \SI{30}{\volt}). Се специфицира
отпорник со $\ge\SI{1}{\watt}$ за резерва.

\subsection{Избор на компоненти}
$C_s$ мора да биде со C0G/NP0 диелектрик: Class-2 (X7R) керамиките губат поголем дел
од својата капацитивност под DC bias и самозагреваат под високофреквентното бранување,
додека C0G е без загуби и стабилен. Компонента со \SI{2.2}{\nano\farad}, \SI{100}{\volt}
C0G е основната (компонента од \SI{200}{\volt} дава дополнителна резерва против
пречекорувањето од ѕвонење). $R_s$ мора да биде неиндуктивен: SMD thin/thick-film чип
(1206 или 2512), никогаш wirewound, така што неговата сопствена индуктивност да не го
поништи snubber-от~\cite{severns}. RC се поставува директно паралелно со drain--source
на секој MOSFET со најкратка можна јамка. Според спецификацијата, сите шест RC места се
поставени како \textbf{do-not-populate (DNP)} footprint-и: ако прототипот покаже
прифатливо ѕвонење тие остануваат празни, а ако не, се монтираат користејќи ги
вредностите измерени на тест-маса од~\eqref{eq:extract}--\eqref{eq:rs}.
Сл.~\ref{fig:snub} ги прикажува поставувањето и ефектот.

\begin{figure}[t]
\centering
\begin{circuitikz}[scale=0.95,transform shape,american]
% low-side device as box
\draw (0,3.4) node[anchor=east]{склопен јазол} -- (1.2,3.4);
\draw (1.2,2.6) rectangle (2.2,3.4);
\draw (1.7,3.0) node{\small $Q$};
\draw (1.7,2.6) -- (1.7,1.6) node[ground]{};
\draw (1.7,3.4) -- (1.7,3.4);
% snubber across D-S
\draw (1.7,3.4) -- (3.0,3.4) to[R,l=$R_s$\,\SI{2.2}{\ohm}] (3.0,2.6)
      to[C,l=$C_s$\,\SI{2.2}{\nano\farad}] (3.0,1.6) node[ground]{};
\draw (1.7,3.4) node[circ]{};
\draw (3.0,3.4) node[anchor=south]{\scriptsize DNP RC};
% ring sketch
\begin{scope}[shift={(4.6,1.7)}]
\draw[->] (0,0) -- (3.2,0) node[anchor=north]{\scriptsize $t$};
\draw[->] (0,-0.9) -- (0,1.5) node[anchor=east]{\scriptsize $v$};
\draw[blue,thick] plot coordinates {(0.00,0.85) (0.05,0.73) (0.11,0.46) (0.16,0.10) (0.21,-0.26) (0.27,-0.54) (0.32,-0.68) (0.38,-0.66) (0.43,-0.50) (0.48,-0.24) (0.54,0.06) (0.59,0.33) (0.64,0.50) (0.70,0.56) (0.75,0.49) (0.80,0.31) (0.86,0.08) (0.91,-0.16) (0.96,-0.34) (1.02,-0.44) (1.07,-0.44) (1.12,-0.34) (1.18,-0.17) (1.23,0.03) (1.29,0.20) (1.34,0.33) (1.39,0.37) (1.45,0.33) (1.50,0.22) (1.55,0.06) (1.61,-0.09) (1.66,-0.22) (1.71,-0.29) (1.77,-0.29) (1.82,-0.23) (1.88,-0.12) (1.93,0.01) (1.98,0.13) (2.04,0.21) (2.09,0.24) (2.14,0.22) (2.20,0.15) (2.25,0.05) (2.30,-0.05) (2.36,-0.14) (2.41,-0.19) (2.46,-0.19) (2.52,-0.15) (2.57,-0.08) (2.62,0.00) (2.68,0.08) (2.73,0.14) (2.79,0.16) (2.84,0.15) (2.89,0.10) (2.95,0.04) (3.00,-0.03)};
\draw[red,thick] plot coordinates {(0.00,0.85) (0.05,0.67) (0.11,0.38) (0.16,0.08) (0.21,-0.18) (0.27,-0.35) (0.32,-0.41) (0.38,-0.36) (0.43,-0.25) (0.48,-0.11) (0.54,0.03) (0.59,0.13) (0.64,0.18) (0.70,0.18) (0.75,0.15) (0.80,0.09) (0.86,0.02) (0.91,-0.04) (0.96,-0.07) (1.02,-0.09) (1.07,-0.08) (1.12,-0.06) (1.18,-0.03) (1.23,0.00) (1.29,0.03) (1.34,0.04) (1.39,0.04) (1.45,0.03) (1.50,0.02) (1.55,0.01) (1.61,-0.01) (1.66,-0.02) (1.71,-0.02) (1.77,-0.02) (1.82,-0.01) (1.88,-0.01) (1.93,0.00) (2.04,0.01) (2.20,0.00) (2.62,0.00) (3.00,0.00)};
\node[blue,anchor=west] at (1.3,1.15) {\scriptsize без snubber};
\node[red,anchor=west] at (1.3,0.62) {\scriptsize со snubber};
\end{scope}
\end{circuitikz}
\caption{DNP RC snubber паралелно со секој прекинувач и квалитативниот ефект врз
ѕвонењето на склопниот јазол: придушувачкиот отпорник $R_s=\sqrt{L_p/C_o}$ ја
отстранува резонантната енергија. Прикажаните вредности се први проценки што треба да
се заменат со измерени на тест-маса преку~\eqref{eq:extract}.}
\label{fig:snub}
\end{figure}

\section{Стратегија за шум и филтрирање}\label{sec:noise}
Филтрирањето се применува на секоја точка каде што се генерира чувствителен аналоген
сигнал или каде што се инјектира склопен шум:
\begin{itemize}[leftmargin=1.3em]
\item \textbf{Излез на CSA} (Дел~\ref{sec:csafilter}): лесен RC од \SI{723}{\kilo\hertz}
      (\SI{100}{\ohm}/\SI{2.2}{\nano\farad}) по фаза, потпирајќи се на синхроно
      семплирање наместо на anti-alias филтер што внесува доцнење~\cite{an5346}.
\item \textbf{Влез на CSA}: сите влезни сериски отпорници мора да бидат \emph{совпаднати}
      меѓу IN+ и IN-- за да се зачува CMRR; enhanced PWM rejection на INA240 е
      примарната одбрана против CM преодната појава, па тежок влезен филтер намерно се
      избегнува~\cite{ina240ds}.
\item \textbf{Разделник на сабирницата}: RC од \SI{1.6}{\kilo\hertz} (Дел~\ref{sec:vfilter}).
\item \textbf{Фазни разделници}: RC од 2--\SI{3}{\kilo\hertz} со компензација на
      фазното доцнење во firmware, плюс Schottky ограничувачи.
\item \textbf{Аналогно напојување}: ferrite bead меѓу $V_{\mathrm{DD}}$ и
      $V_{\mathrm{DDA}}$ со декуплирање \SI{1}{\micro\farad}\,$+$\,\SI{100}{\nano\farad}
      на $V_{\mathrm{DDA}}$; посебно \SI{1}{\micro\farad}\,$+$\,\SI{100}{\nano\farad} на
      $V_{\mathrm{REF+}}$, и декуплирањето на VREFBUF специфицирано од ST ако се
      користи внатрешната референца~\cite{stm32g473ds}.
\item \textbf{Напојување на INA240}: \SI{100}{\nano\farad}\,$+$\,\SI{1}{\micro\farad}
      поставени на приклучокот.
\item \textbf{Декуплирање на DC-сабирницата}: керамики со ниска ESL паралелно со bulk
      електролитите, блиску до секој half-bridge, за да се минимизира индуктивноста на
      комутационата јамка $L_p$ --- ова го напаѓа \emph{изворот} на ѕвонењето што
      snubber-от инаку мора да го придушува~\cite{severns}.
\item \textbf{Заземјување/layout}: единствен (star) спој меѓу аналогната земја (AGND)
      и енергетската земја (PGND); Kelvin мерење на shunt-от; кратки, тесно спрегнати
      диференцијални сигнални траги; и компактна склопна јамка. Овие мерки во layout-от
      се исто толку важни колку и филтрите со компоненти за инвертор со
      \SI{40}{\kilo\hertz} и \SI{>100}{\ampere}.
\end{itemize}

\section{Спецификација на материјали (BOM)}\label{sec:bom}
Спецификацијата на материјали се состои од два дела што се надополнуваат: пасивните
помошни компоненти за управување на гејтот, веќе наведени по гранка и по плоча во
Табела~\ref{tab:gatebom}, и компонентите на аналогниот влезен степен со LCSC
нарачкини броеви, дадени во Табела~\ref{tab:bom}. Активните уреди се
\textbf{$12\times$ IPT015N10N5} (четири по гранка, нарачкин број LCSC C108964) и
\textbf{$3\times$ UCC27211A-Q1} (по еден на фаза), заедно со \textbf{$3\times$ INA240A2}
(C129949) за мерење на струја.

Активните уреди (INA240A2PWR, IPT015N10N5) и стандардните чип отпорници/кондензатори
се производи со голем обем што рутински имаат залиха далеку над \num{500} единици;
единствената компонента што треба да се потврди при пуштање на нарачката е shunt-от со
голема моќност, за кој е дадена алтернатива составена од два крајно вообичаени
\SI{0.2}{\milli\ohm} \(2512\) елементи паралелно ($=\SI{0.1}{\milli\ohm}$,
$\approx\SI{0.93}{\watt}$ по секој). Количините претпоставуваат една трифазна плоча;
snubber компонентите се DNP (само footprint-и).

\begin{table*}[t]
\centering
\caption{Спецификација на материјали за аналогниот влезен степен. Залихата се менува; потврди при пуштање на нарачката.}
\label{tab:bom}
\small
\renewcommand{\arraystretch}{1.2}
\begin{tabular}{p{3.1cm} p{1.0cm} p{3.2cm} p{2.6cm} p{2.0cm} p{2.7cm}}
\toprule
\textbf{Функција} & \textbf{Кол.} & \textbf{Компонента / вредност} & \textbf{Клучен номинал} & \textbf{LCSC} & \textbf{Забелешка}\\
\midrule
Струен засилувач & 3 & TI INA240A2PWR & $G{=}50$, $-4..\SI{80}{\volt}$ CM & C129949 & по еден на фаза~\cite{lcsc_ina}\\
Фазен shunt & 3 & \SI{0.1}{\milli\ohm} метален елемент, \SI{5}{\watt} & $\le\SI{2}{\nano\henry}$ ESL, Kelvin & потврди & на пр.\ Milliohm 6.6$\times$6.6\,mm; види алт.\ подолу~\cite{lcsc_shunt}\\
\;\;(алтернатива за shunt) & 6 & \SI{0.2}{\milli\ohm} \(2512\) $\times2$ \(\parallel\) & \SI{2}{\watt} по сек. & голема залиха & крајно вообичаени; $=\SI{0.1}{\milli\ohm}$\\
Излезен филтер R на CSA & 3 & \SI{100}{\ohm} \(0603\) 1\% & --- & голема залиха & со $C_f$\\
Излезен филтер C на CSA & 3 & \SI{2.2}{\nano\farad} C0G \(0603\) & \SI{50}{\volt} & голема залиха & $f_c\,\SI{723}{\kilo\hertz}$\\
Горен дел на разделник & 4(8) & \SI{110}{\kilo\ohm} (2$\times$\SI{55}{\kilo\ohm}) \(0805\) 1\% & $\ge\SI{150}{\volt}$ & голема залиха & сабирница + 3 фази\\
Долен дел на разделник & 4 & \SI{4.7}{\kilo\ohm} \(0603\) 1\% & --- & голема залиха & сооднос $24.4$\\
Филтер кондензатор на сабирница & 1 & \SI{22}{\nano\farad} C0G/X7R \(0603\) & \SI{50}{\volt} & голема залиха & $f_c\,\SI{1.6}{\kilo\hertz}$\\
Филтер кондензатор на фаза & 3 & \SI{22}{\nano\farad} \(0603\) & \SI{50}{\volt} & голема залиха & $f_c\approx\SI{2.5}{\kilo\hertz}$\\
Ограничувачка диода на ADC & 4 & BAT54S двојна Schottky & \SI{30}{\volt} & голема залиха & кон $V_{\mathrm{DDA}}$/GND\\
Snubber R (DNP) & 6 & \SI{2.2}{\ohm} \(1206\) неинд., $\ge\SI{1}{\watt}$ & --- & голема залиха & само footprint\\
Snubber C (DNP) & 6 & \SI{2.2}{\nano\farad} C0G \(1206\) & \SI{100}{\volt} (\SI{200}{\volt} опц.) & голема залиха & само footprint\\
Ferrite bead на VDDA & 1 & \SI{600}{\ohm}@\SI{100}{\mega\hertz} \(0805\) & --- & голема залиха & $V_{\mathrm{DD}}\!\to\!V_{\mathrm{DDA}}$\\
Сет за декуплирање & --- & \SI{100}{\nano\farad}/\SI{1}{\micro\farad} \(0402/0603\) & \SI{16}{\volt}+ & голема залиха & CSA, VDDA, VREF\\
Енергетски MOSFET (реф.) & 12 & Infineon IPT015N10N5 & \SI{100}{\volt}/\SI{300}{\ampere}/\SI{1.5}{\milli\ohm} & C108964 & енергетски степен, 4 по гранка~\cite{lcsc_ipt}\\
\bottomrule
\end{tabular}
\end{table*}

\begin{thebibliography}{99}
\bibitem{stm32g473ds} STMicroelectronics, ``STM32G473xB/xC/xE datasheet (DS12712),'' ревизија~5, 2023.
[Онлајн]. Достапно на: \url{https://www.st.com/resource/en/datasheet/stm32g473cb.pdf}

\bibitem{ucc27211} Texas Instruments, ``UCC27211A-Q1 Automotive 120-V, 3.7-A/4.5-A half-bridge
driver with 8-V UVLO,'' datasheet SLUSCG0C, Dec.\,2015 --- rev.\,Sep.\,2025. [Онлајн]. Достапно на:
\url{https://www.ti.com/lit/ds/symlink/ucc27211a-q1.pdf}

\bibitem{ipt015ds} Infineon Technologies, ``IPT015N10N5 OptiMOS\texttrademark~5 \SI{100}{\volt}
power transistor,'' final data sheet, ревизија~2.4, 2023. [Онлајн]. Достапно на:
\url{https://www.infineon.com/assets/row/public/documents/24/49/infineon-ipt015n10n5-datasheet-en.pdf}

\bibitem{ipt015part} Infineon Technologies, ``IPT015N10N5 product page'' (страница на производот,
параметарски преглед; \(C_{\mathrm{oss}}=\SI{1800}{\pico\farad}\), \SI{100}{\volt}, \SI{300}{\ampere}).
[Онлајн]. Достапно на: \url{https://www.infineon.com/part/IPT015N10N5}

\bibitem{onsemi_bldc} onsemi, ``Trapezoidal control of BLDC motors,'' технички блог. [Онлајн].
Достапно на: \url{https://www.onsemi.com/company/news-media/blog/industrial/en-us/trapezoidal-control-of-bldc-motors}

\bibitem{nxp_backemf} NXP Semiconductors, ``3-phase BLDC motor control with sensorless back-EMF
zero-crossing detection,'' апликациона белешка AN1913. [Онлајн]. Достапно на:
\url{https://www.nxp.com/docs/en/application-note/AN1913.pdf}

\bibitem{balogh} L.~Balogh, ``Fundamentals of MOSFET and IGBT gate driver circuits,'' Texas
Instruments application report SLUA618A (ревизија на SLUP169), 2017--2018. [Онлајн]. Достапно на:
\url{https://www.ti.com/lit/ml/slua618/slua618.pdf}

\bibitem{infineon_monolithic} Infineon Technologies, ``Using monolithic high-voltage gate drivers''
(димензионирање на bootstrap, отпорници на гејтот и насоки за полумост; наследник на IR белешките
AN-978/AN-1123), апликациона белешка. [Онлајн]. Достапно на:
\url{https://www.infineon.com/assets/row/public/documents/24/42/infineon-using-monolithic-high-voltage-gate-drivers-applicationnotes-en.pdf}

\bibitem{infineon_bootoverview} Infineon Technologies, ``Bootstrap circuits --- overview and
selection guide,'' Knowledge Base статија. [Онлајн]. Достапно на:
\url{https://community.infineon.com/t5/Knowledge-Base-Articles/Bootstrap-circuits-overview-and-selection-guide/ta-p/720580}

\bibitem{ucc27211folder} Texas Instruments, ``UCC27211A product folder.'' [Онлајн]. Достапно на:
\url{https://www.ti.com/product/UCC27211A}

\bibitem{infineon_bootcalc} Infineon Technologies, ``A typical calculation of a bootstrap capacitor's
value,'' Knowledge Base статија. [Онлајн]. Достапно на:
\url{https://community.infineon.com/t5/Knowledge-Base-Articles/A-typical-calculation-of-a-bootstrap-capacitor-s-value/ta-p/739554}

\bibitem{infineon_maxboot} Infineon Technologies, ``Calculating maximum bootstrap capacitance for
integrated-diode gate drivers,'' Knowledge Base статија. [Онлајн]. Достапно на:
\url{https://community.infineon.com/t5/Knowledge-Base-Articles/Calculating-maximum-bootstrap-capacitance-for-integrated-diode-gate-drivers/ta-p/1015593}

\bibitem{pmeg4030er} Nexperia, ``PMEG4030ER --- \SI{40}{\volt}, \SI{3}{\ampere} low-$V_F$ Schottky
barrier rectifier,'' технички лист, 2023. [Онлајн]. Достапно на:
\url{https://assets.nexperia.com/documents/data-sheet/PMEG4030ER.pdf}

\bibitem{nexperia_gdfund} Nexperia, ``Power MOSFET gate driver fundamentals,'' апликациона белешка
AN90059, 2025. [Онлајн]. Достапно на:
\url{https://assets.nexperia.com/documents/application-note/AN90059.pdf}

\bibitem{toshiba_dvdt} Toshiba Electronic Devices \& Storage, ``Impacts of the dv/dt rate on
MOSFETs,'' апликациона белешка, 2018. [Онлајн]. Достапно на:
\url{https://toshiba.semicon-storage.com/info/docget.jsp?did=59464}

\bibitem{ti_lowside_threephase} Texas Instruments, ``Low-drift, low-side current measurements for
three-phase systems,'' апликациона белешка (од документацијата на INA240). [Онлајн]. Достапно на:
\url{https://www.ti.com/product/INA240}

\bibitem{tida00913} Texas Instruments, ``TIDA-00913: 48-V/10-A three-phase inverter with in-line
shunt-based phase current sensing reference design.'' [Онлајн]. Достапно на:
\url{https://www.ti.com/tool/TIDA-00913}

\bibitem{ina240ds} Texas Instruments, ``INA240 $-4$\,V to $80$\,V, bidirectional, ultra-precise
current-sense amplifier with enhanced PWM rejection (SBOS662),'' технички лист. [Онлајн]. Достапно на:
\url{https://www.ti.com/lit/gpn/INA240}

\bibitem{an2834} STMicroelectronics, ``AN2834: How to get the best ADC accuracy in STM32
microcontrollers.'' [Онлајн]. Достапно на:
\url{https://www.st.com/resource/en/application_note/an2834-how-to-get-the-best-adc-accuracy-in-stm32-microcontrollers-stmicroelectronics.pdf}

\bibitem{an5346} STMicroelectronics, ``AN5346: STM32G4 ADC use tips and recommendations.'' [Онлајн].
Достапно на:
\url{https://www.st.com/resource/en/application_note/an5346-stm32g4-adc-use-tips-and-recommendations-stmicroelectronics.pdf}

\bibitem{an5306} STMicroelectronics, ``AN5306: Operational amplifier (OPAMP) usage in STM32G4 Series.''
[Онлајн]. Достапно на:
\url{https://www.st.com/resource/en/application_note/an5306-operational-amplifier-opamp-usage-in-stm32g4-series-stmicroelectronics.pdf}

\bibitem{severns} R.~Severns, ``Design of snubbers for power circuits,'' Cornell Dubilier апликациона
белешка. [Онлајн]. Достапно на: \url{https://www.cde.com/resources/technical-papers/design.pdf}

\bibitem{digikey_snubber} DigiKey, ``Resistor-capacitor (RC) snubber design for power switches,''
технички напис. [Онлајн]. Достапно на:
\url{https://www.digikey.com/en/articles/resistor-capacitor-rc-snubber-design-for-power-switches}

\bibitem{lcschome} LCSC Electronics, почетна страница на дистрибутер (\num{560000}+ SKU, \(2600+\)
производители). [Онлајн]. Достапно на: \url{https://www.lcsc.com/}

\bibitem{lcsc_ina} LCSC Electronics, ``INA240A2PWR (C129949).'' [Онлајн]. Достапно на:
\url{https://www.lcsc.com/product-detail/C129949.html}

\bibitem{lcsc_ipt} LCSC Electronics, ``IPT015N10N5 (C108964).'' [Онлајн]. Достапно на:
\url{https://www.lcsc.com/product-detail/C108964.html}

\bibitem{lcsc_shunt} LCSC Electronics, ``Категорија: current sense / shunt отпорници.'' [Онлајн].
Достапно на: \url{https://www.lcsc.com/products/Current-Sense-Resistors-Shunt-Resistors_909.html}
\end{thebibliography}

\end{document}
